% Options for packages loaded elsewhere
\PassOptionsToPackage{unicode}{hyperref}
\PassOptionsToPackage{hyphens}{url}
\documentclass[
]{article}
\usepackage{xcolor}
\usepackage{amsmath,amssymb}
\setcounter{secnumdepth}{-\maxdimen} % remove section numbering
\usepackage{iftex}
\ifPDFTeX
  \usepackage[T1]{fontenc}
  \usepackage[utf8]{inputenc}
  \usepackage{textcomp} % provide euro and other symbols
\else % if luatex or xetex
  \usepackage{unicode-math} % this also loads fontspec
  \defaultfontfeatures{Scale=MatchLowercase}
  \defaultfontfeatures[\rmfamily]{Ligatures=TeX,Scale=1}
\fi
\usepackage{lmodern}
\ifPDFTeX\else
  % xetex/luatex font selection
\fi
% Use upquote if available, for straight quotes in verbatim environments
\IfFileExists{upquote.sty}{\usepackage{upquote}}{}
\IfFileExists{microtype.sty}{% use microtype if available
  \usepackage[]{microtype}
  \UseMicrotypeSet[protrusion]{basicmath} % disable protrusion for tt fonts
}{}
\makeatletter
\@ifundefined{KOMAClassName}{% if non-KOMA class
  \IfFileExists{parskip.sty}{%
    \usepackage{parskip}
  }{% else
    \setlength{\parindent}{0pt}
    \setlength{\parskip}{6pt plus 2pt minus 1pt}}
}{% if KOMA class
  \KOMAoptions{parskip=half}}
\makeatother
\usepackage{longtable,booktabs,array}
\newcounter{none} % for unnumbered tables
\usepackage{calc} % for calculating minipage widths
% Correct order of tables after \paragraph or \subparagraph
\usepackage{etoolbox}
\makeatletter
\patchcmd\longtable{\par}{\if@noskipsec\mbox{}\fi\par}{}{}
\makeatother
% Allow footnotes in longtable head/foot
\IfFileExists{footnotehyper.sty}{\usepackage{footnotehyper}}{\usepackage{footnote}}
\makesavenoteenv{longtable}
\setlength{\emergencystretch}{3em} % prevent overfull lines
\providecommand{\tightlist}{%
  \setlength{\itemsep}{0pt}\setlength{\parskip}{0pt}}
\usepackage{bookmark}
\IfFileExists{xurl.sty}{\usepackage{xurl}}{} % add URL line breaks if available
\urlstyle{same}
\hypersetup{
  hidelinks,
  pdfcreator={LaTeX via pandoc}}

\author{}
\date{}

\begin{document}

\section{Empirical Boundaries of Post-Training Quantization and
Throughput Saturation for ViT-L on Blackwell GPUs: A Three-Tool-Chain
Convergence Study with TensorRT
10.13}\label{empirical-boundaries-of-post-training-quantization-and-throughput-saturation-for-vit-l-on-blackwell-gpus-a-three-tool-chain-convergence-study-with-tensorrt-1013}

\begin{quote}
\textbf{Single-file paper draft V1.0.1 (2026-05-04)} --- extends V1.0.0
(2026-05-01) with new §4.7 Throughput-Oriented Serving and Pool
Saturation, integrating V1.0.3 empirical findings (5.24× total speedup
vs FP32, BF16 dense path SM 96--99\% saturation across all tested
regimes, 720 qps G1 in-process pool ceiling, perfect bit-exact
concurrent execution under N=2/N=4 load). Abstract bilingual updated
with V1.0.3 throughput results; §5.3 V1.3 QAT motivation strengthened
with utilization-saturation evidence; conclusion §7 reflects V1.0.3
closure status (5 of 7 V1.0.3 SMART goals fully passed, 1 close at
hardware ceiling, 2 physics-blocked, 1 user-blocker).

\textbf{Authors}: {[}Author Name{]}, PolyU

\textbf{Date}: 2026-05

\textbf{Source drafts} (in \texttt{Wiki/2-技术报告/}):

\begin{itemize}
\tightlist
\item
  \texttt{paper\_abstract\_intro\_draft\_V1.0.0.md} --- Abstract (EN +
  中文) + § 1 Introduction
\item
  \texttt{paper\_literature\_review\_draft\_V1.0.0.md} --- § 2
  Literature Review
\item
  \texttt{paper\_methodology\_draft\_V1.0.0.md} --- § 3 Methodology
\item
  \texttt{paper\_results\_draft\_V1.0.0.md} --- § 4 Results
\item
  \texttt{paper\_discussion\_draft\_V1.0.0.md} --- § 5 Discussion
\item
  \texttt{paper\_limitations\_conclusion\_draft\_V1.0.0.md} --- § 6
  Limitations + § 7 Conclusion
\end{itemize}

\textbf{Figure references} (paths relative to
\texttt{Code/Artifacts/reports/figures/}):

\begin{itemize}
\tightlist
\item
  Figure 1: \texttt{benchmark\_trtexec\_bf16\_speedup.svg}
  (multi-resolution speedup bars)
\item
  Figure 2: \texttt{benchmark\_cpp\_runtime\_speedup.svg} (C++
  end-to-end speedup)
\item
  Figure 3: \texttt{benchmark\_bf16\_cosine\_min.svg} (multi-resolution
  cosine min)
\item
  Figure 4: \texttt{benchmark\_bf16\_cosine\_mean.svg} (multi-resolution
  cosine mean)
\item
  Figure 5: \texttt{benchmark\_bf16\_vs\_int8\_tradeoff.svg} (12-point
  INT8 sensitivity scatter)
\item
  Figure 6: \texttt{layer\_ablation\_diversity\_vs\_balance.svg}
  (4-layer ablation)
\end{itemize}

\textbf{Reproducibility kit}:

\begin{itemize}
\tightlist
\item
  87-row machine-readable benchmark matrix:
  \texttt{Code/Artifacts/reports/formal\_benchmark\_matrix.csv}
\item
  Atomic SHA256 manifest covering 438+ files:
  \texttt{Code/Artifacts/reports/artifact\_manifest\_formal\_with\_sha256.json}
\item
  336 pytest passing + 81\% line coverage across 114 source files
\item
  One-shot Windows pipeline:
  \texttt{Code/scripts/run\_formal\_hf\_pipeline\_windows.ps1}
\end{itemize}
\end{quote}

\begin{center}\rule{0.5\linewidth}{0.5pt}\end{center}

\subsection{English Abstract}\label{english-abstract}

\subsubsection{Empirical Boundaries of Post-Training Quantization for
ViT-L on Blackwell GPUs: A Three-Tool-Chain Convergence Study with
TensorRT
10.13}\label{empirical-boundaries-of-post-training-quantization-for-vit-l-on-blackwell-gpus-a-three-tool-chain-convergence-study-with-tensorrt-1013}

\textbf{Background.} Visual self-supervised foundation models such as
DINOv3 ViT-L/16 LVD-1689M expose multi-scale intermediate features for
dense prediction downstream tasks (depth, segmentation), but FP32
inference cost on consumer-grade hardware prohibits deployment.

\textbf{Purpose.} We empirically map the post-training quantization
(PTQ) boundaries of ViT-L/16 on NVIDIA Blackwell sm\_120 with TensorRT
10.13, and establish whether mixed-precision strategies can overcome
upstream quantization noise to satisfy a stringent (cosine ≥ 0.99 ∧
speedup ≥ 2.2×) deployment target.

\textbf{Method.} We benchmark 12 inference precision candidates---FP32,
BF16, FP16, FP8 default and partial, INT8 default and three node-level
partial variants, INT8 SmoothQuant α-sweep at three settings, and three
independent mixed-precision strategies (PyTorch ModelOpt
\texttt{disable\_quantizer}, TensorRT \texttt{-\/-layerPrecisions}, and
ONNX-level Q/DQ stripping)---across three resolutions (224, 336, 518) on
RTX 5080. Cosine similarity against the FP32 baseline is measured on
1,000 real images per (resolution, candidate) pair. Cross-language
Python ↔ C++ runtime parity is verified on bit-identical engine binaries
with deterministic input.

\textbf{Findings.} BF16 prefer is the only candidate inside the (cos ≥
0.99 ∧ speedup ≥ 2.2×) target region, achieving 3.86× peak speedup at
r518 batch 8 with feat\_layer\_20 cosine ≥ 0.998 across all three
resolutions. The three independent mixed-precision tool chains produce
numerically equivalent results (cos\_min within 0.0005 across paths, b8
latency within 0.02× when both force fp32 fallback), demonstrating that
the precision bottleneck is upstream cumulative INT8 quantization noise
from blocks 0--15, not the choice of mixed-precision tool. Python ↔ C++
runtime parity holds bit-identically across all three resolutions at
batch 1.

\textbf{Throughput Extension (V1.0.3).} Beyond single-stream latency, we
extend the BF16 baseline with a multi-context inference pool
(\texttt{TRTInfererPool}, ADR-020 Phase 2: shared engine + N independent
execution contexts + per-slot CUDA streams) and characterize aggregate
throughput plus GPU utilization at four operating points. End-to-end
aggregate throughput reaches \textbf{720 qps at r224 b8 N=4 slots × 16
caller threads} (2.094× over a V1.0.2 single-stream baseline of 343 qps;
5.24× over the original FP32 trtexec gpu-compute baseline of 137 qps).
GPU utilization measurements at 100 ms cadence confirm that the pool
drives the SM occupancy to \textbf{98.74\% mean / 100\% p95} at this
operating point --- the in-process pool is at the hardware compute
ceiling, and the residual \textasciitilde80 qps gap to the original 800
qps target plan threshold is bounded by host-device transfer overhead
rather than GPU compute. Concurrent execution under N=2 / N=4 produces
output that is bit-identical (max\_abs\_error = 0, cosine = 1.0) to the
single-context reference. Three of three tested resolutions (224, 336,
518) operate at SM 88--99\% under their respective optimal batch points;
r336 and r518 saturate at SM 96--99\% even at N=1, leaving no headroom
for multi-context to exploit.

\textbf{Implications.} Mixed-precision PTQ is fundamentally insufficient
for tight cosine constraints on this model--hardware combination; the
V1.0.3 throughput results extend this conclusion to the runtime axis,
showing that the BF16 dense path is also fully saturated at the GPU
compute level. Quantization-aware fine-tuning is therefore the only
remaining path on either axis, and we publish a complete V1.3 design ADR
with a four-condition launch threshold. We also release a 56-row
machine-readable benchmark matrix, twelve V1.0.3 throughput data points,
eight figures, 271+ unit tests, and an atomic SHA256 manifest covering
419+ artifact files for community reproducibility.

\textbf{Keywords}: TensorRT inference acceleration · Vision Transformer
ViT-L · INT8 post-training quantization · mixed-precision quantization ·
DINOv3 self-supervised vision · Blackwell GPU architecture · DPT-style
multi-scale fusion

\emph{Word count: 480 words (V1.0.0: 312 + V1.0.3 extension: 168)}

\subsection{简体中文摘要（独立撰写，非英文翻译）}\label{ux7b80ux4f53ux4e2dux6587ux6458ux8981ux72ecux7acbux64b0ux5199ux975eux82f1ux6587ux7ffbux8bd1}

\subsubsection{Blackwell GPU 上 ViT-L 推理后量化的实证边界：基于
TensorRT 10.13
的三工具链等价性研究}\label{blackwell-gpu-ux4e0a-vit-l-ux63a8ux7406ux540eux91cfux5316ux7684ux5b9eux8bc1ux8fb9ux754cux57faux4e8e-tensorrt-1013-ux7684ux4e09ux5de5ux5177ux94feux7b49ux4ef7ux6027ux7814ux7a76}

\textbf{研究背景。} DINOv3 ViT-L/16 LVD-1689M
等视觉自监督基础模型通过暴露多尺度中间特征支持深度估计、语义分割等下游密集预测任务，但
FP32 精度的推理成本使其难以在消费级硬件上部署。

\textbf{研究目的。} 在 NVIDIA Blackwell sm\_120 架构 GPU 配合 TensorRT
10.13 推理引擎的工程组合下，本研究系统刻画 ViT-L/16
推理后量化（PTQ）的实证操作边界，并实证检验混合精度策略能否突破前段累积量化噪声、满足"cosine
≥ 0.99 同时加速 ≥ 2.2×"的严格部署阈值。

\textbf{研究方法。} 在 RTX 5080 上对 12 个精度候选（FP32 / BF16 / FP16 /
FP8 默认与节点级 / INT8 默认与三种节点级 / INT8 SmoothQuant 三档 alpha /
三种独立工具链的混合精度策略）进行多分辨率（224 / 336 /
518）系统化基准测试。每个（分辨率，候选）组合在 1000 张真实图片上对照
FP32 基线计算 4 个中间层输出的余弦相似度，并通过同一 engine 二进制 +
确定性输入验证 Python 与 C++ 推理运行时的位级一致性。

\textbf{研究发现。} BF16 prefer 是唯一进入"cos ≥ 0.99 同时加速 ≥
2.2×"目标区域的候选，在 r518 batch 8 取得 3.86× 加速峰值，且三档分辨率下
feat\_layer\_20 余弦相似度均 ≥
0.998。三种相互独立的混合精度工具链（PyTorch ModelOpt 的
\texttt{disable\_quantizer}、TensorRT 命令行
\texttt{-\/-layerPrecisions}、ONNX 图层 Q/DQ
节点剥离）产出数值等价的结果------三路径 cos\_min 差异 \textless{}
0.0005、batch 8 延迟差异 \textless{} 0.02× --- 实证表明精度瓶颈在于
blocks 0-15 累积的 INT8 量化噪声，与所选用的混合精度工具无关。Python 与
C++ 运行时在三档分辨率下 batch 1 全部位级一致。

\textbf{吞吐量扩展研究（V1.0.3）。} 在 BF16
单流加速基础上，本研究进一步扩展多 context
推理池架构（\texttt{TRTInfererPool}，ADR-020 Phase 2：共享 engine + N
个独立 execution context + 每槽位独立 CUDA
stream），并在四个工作点上系统刻画 aggregate 吞吐量与 GPU
利用率。端到端聚合吞吐量在 \textbf{r224 b8 N=4 slots × 16 caller
threads} 配置下达到 \textbf{720 qps}（相比 V1.0.2 单流基线 343 qps 提升
2.094×；相比原始 FP32 trtexec gpu-compute 基线 137 qps 提升 5.24×）。100
ms 采样间隔的 GPU 利用率监测确认池化架构将 SM 占用率推至 \textbf{平均
98.74\% / p95 100\%} ------ 在该工作点上 in-process
池已达硬件计算上限，与原计划阈值 800 qps 之间残余 \textasciitilde80 qps
的差距由 host-device 数据传输开销而非 GPU 计算决定。N=2 / N=4
并发执行的输出与单 context 参考严格位级一致（max\_abs\_error = 0、cosine
= 1.0）。三档分辨率（224 / 336 / 518）在各自最优批量下均运行于 SM
88--99\% 区间；r336 和 r518 在 N=1 时 SM 已达 96--99\%，多 context
无空闲计算可填补。

\textbf{研究意义。} 混合精度 PTQ
在该模型-硬件组合下根本上无法满足严格余弦约束；V1.0.3
吞吐量结果在运行时维度上进一步延伸该结论 ------ BF16 dense path 在 GPU
计算层面同样达到完全饱和。量化感知微调（QAT）因此在精度与运行时两个维度上均成为唯一剩余路径。本研究公开完整的
V1.3 设计 ADR，包含四条启动门槛。我们同时发布 56 行机器可读 benchmark
矩阵、12 个 V1.0.3 吞吐量数据点、8 张可视化产物、271+ 单元测试、覆盖
419+ 文件的原子化 SHA256 完整索引，供学术社区复现使用。

\textbf{关键词}：TensorRT 推理加速 · Vision Transformer ViT-L · INT8
后量化 · 混合精度量化 · DINOv3 自监督视觉模型 · Blackwell GPU 架构 ·
DPT-style 多尺度融合 · 多 context 推理池 · GPU 利用率饱和

\emph{字数：约 1080 字（V1.0.0：720 字 + V1.0.3 扩展：360 字）}

\subsection{1. Introduction（Actual Draft, \textasciitilde1500 words
English）}\label{1-introductionactual-draft-1500-words-english}

\subsubsection{1.1 Context and
Background}\label{11-context-and-background}

The deployment of large self-supervised vision foundation models on
consumer-grade GPUs has emerged as a critical engineering bottleneck for
production computer vision systems. Models such as DINOv3 ViT-L/16
LVD-1689M {[}Meta AI, 2024{]}, pretrained on the Large Visual Dataset of
1.689 billion images, provide rich multi-scale intermediate
representations that downstream dense prediction architectures (DPT
{[}Ranftl et al., 2021{]}, depth estimation, semantic segmentation)
consume directly via four intermediate-layer hooks at transformer blocks
4, 12, 16, and 20. The 24-block ViT-Large backbone with its
1024-dimensional hidden states and 197-token spatial contract (after
register-token trimming) presents simultaneously high inference cost and
high deployment value: a single forward pass at 224×224 input batch 8
takes 28 ms on FP32 on RTX 5080, while the four-output multi-scale
features it produces enable dense prediction tasks that single-output
classification backbones cannot serve.

NVIDIA\textquotesingle s TensorRT inference compiler {[}NVIDIA
Corporation, 2026{]} has been the de facto path for production-grade
acceleration of such models, with the 10.x release line introducing
Blackwell sm\_120 hardware support and explicit Q/DQ INT8 quantization
through the Model Optimizer toolkit. The combination promises 2-5×
speedups via reduced-precision execution while preserving the numerical
fidelity required by downstream tasks. However, the engineering reality
at the intersection of three relatively new components---Blackwell
sm\_120 architecture (released 2025), TensorRT 10.13.2.6 (released early
2026), and ViT-L architecture as a target---has not been systematically
mapped in the public literature.

Visual foundation models present an additional precision challenge:
downstream dense prediction quality is highly sensitive to perturbations
in intermediate features. While image classification can tolerate cosine
similarities above 0.95 against the FP32 baseline, depth estimation and
segmentation downstream heads exhibit observable degradation below 0.99
cosine on the deepest hooked layer, motivating a stringent precision
target that standard PTQ literature rarely interrogates.

\subsubsection{1.2 Problem Statement}\label{12-problem-statement}

This study investigates the operating region defined by two simultaneous
constraints: (1) feat\_layer\_20 cosine similarity ≥ 0.99 against the
FP32 baseline on representative real images, and (2) inference latency
speedup ≥ 2.2× on RTX 5080 (Blackwell sm\_120) at TensorRT 10.13.2.6. We
refer to this region as the "G2 ideal region" throughout the paper.

The dual constraint is empirically narrow: the standard PTQ literature
on Vision Transformers reports cosine ≥ 0.95 at speedups in the 2-3×
range {[}Frantar et al., 2023; Xiao et al., 2023; Liu et al., 2024{]},
but cos ≥ 0.99 at the deepest layer of ViT-L/16 with ≥ 2.2× speedup---on
the specific Blackwell + TRT 10.13 + DINOv3 ViT-L combination---has not
been demonstrated.

The problem is sharpened by three observations from preliminary
experiments: (a) FP16 reduced precision, despite passing trtexec sanity
checks, produces all-NaN intermediate features under formal weights on
this hardware-software combination, eliminating it as a low-precision
candidate; (b) default ModelOpt INT8 PTQ collapses to cos\_mean ≈ 0.20
at the deepest layer even with 500-image calibration; and (c) node-level
partial INT8 quantization (limiting INT8 to a single MatMul subgraph)
recovers cosine quality but reduces speedup to 1.04-1.08×, eliminating
its value as a deployment candidate. These observations motivate the
deeper investigation that constitutes this paper.

\subsubsection{1.3 Research Gap}\label{13-research-gap}

Three gaps in the existing literature motivate this work.

\textbf{First}, the public TensorRT benchmark literature for ViT-L on
Blackwell sm\_120 is sparse. The architecture has been available for
less than a year at the time of this work, and the intersection with the
latest TRT 10.13 release contains undocumented behaviors---such as BF16
incompatibility with explicit Q/DQ ONNX in the Myelin Fill kernel---that
affect engineering decisions but are not surfaced in vendor
documentation or community benchmarks.

\textbf{Second}, ViT-L INT8 PTQ studies in the literature predominantly
target cosine ≥ 0.95 thresholds suitable for image classification, using
techniques such as GPTQ {[}Frantar et al., 2023{]}, AWQ, and SmoothQuant
{[}Xiao et al., 2023{]}. These methods have been validated on language
and vision models but the cosine ≥ 0.99 stringent regime on
ViT-L/16\textquotesingle s deepest hook has not been probed
comprehensively.

\textbf{Third}, when standard PTQ fails to meet a target precision,
mixed-precision strategies are typically proposed as a remedy: force the
most error-sensitive transformer blocks back to higher precision while
keeping the remainder INT8. The literature, however, evaluates such
mixed-precision strategies one tool chain at a time. Whether the choice
of tool chain (PyTorch quantizer disable, TensorRT command-line
per-layer override, or ONNX-level graph rewrite) materially affects the
precision-speedup trade-off has not been established empirically. This
gap is consequential: if multiple tool chains converge to identical
results, the convergence is itself evidence that the precision
bottleneck lies elsewhere---and indicates where remediation effort
should and should not be directed.

\subsubsection{1.4 Purpose and Research
Questions}\label{14-purpose-and-research-questions}

The purpose of this study is to \textbf{empirically map the (precision,
speedup) operating region of TensorRT 10.13 inference for DINOv3
ViT-L/16 on Blackwell sm\_120}, and to \textbf{identify the binding
constraint} that prevents any post-training quantization candidate from
entering the G2 ideal region.

The following research questions guide the investigation:

\begin{itemize}
\item
  \textbf{RQ1}: Among 12 standard inference precision
  candidates---including FP32, BF16, FP16, FP8 default and partial, INT8
  default and node-level partial variants, INT8 SmoothQuant at three α
  settings, and three mixed-precision strategies---which fall inside the
  (cos ≥ 0.99 ∧ speedup ≥ 2.2×) target region on this model + hardware
  combination?
\item
  \textbf{RQ2}: Does the choice of tool chain (PyTorch
  ModelOpt\textquotesingle s \texttt{disable\_quantizer}, TensorRT
  command-line \texttt{-\/-layerPrecisions}, or ONNX library-level Q/DQ
  stripping) materially affect the precision-speedup trade-off of
  mixed-precision strategies that force a subset of transformer blocks
  back to higher precision?
\item
  \textbf{RQ3}: For candidates that fail the G2 target, what empirical
  evidence determines the binding constraint, and what implementation
  path could lift it without replacing the underlying ViT-L
  architecture?
\end{itemize}

We answer these questions through a tightly controlled benchmark
protocol with locked GPU clocks, spin-wait latency measurement,
multi-resolution coverage, and 1000-image cosine evaluation. The
protocol\textquotesingle s reproducibility kit---comprising a 56-row
machine-readable benchmark matrix, eight figures, 271 unit tests, and an
atomic SHA256 manifest covering 419+ files---is released alongside this
paper.

\subsubsection{1.5 Significance of the
Study}\label{15-significance-of-the-study}

This study makes three contributions to the inference systems
literature.

\textbf{An empirical map of a new hardware-software intersection.} We
document the (precision, speedup) operating region of TensorRT 10.13 on
Blackwell sm\_120 for DINOv3 ViT-L/16 across three deployment-relevant
resolutions (224, 336, 518). The map identifies BF16 prefer as the only
inference candidate inside the G2 ideal region, with a 3.86× peak
speedup at r518 batch 8 and feat\_layer\_20 cosine ≥ 0.998 throughout.
Counter-intuitively, the deepest hooked layer at the largest resolution
achieves the highest cosine (0.999171), attributable to BF16
quantization noise being diluted across more patch tokens at higher
spatial dimensions.

\textbf{A three-tool-chain equivalence proof for mixed-precision
recovery.} We show, via three independent implementations across PyTorch
ModelOpt, TensorRT command-line, and ONNX library layers, that
mixed-precision strategies forcing transformer blocks 16--19 to non-INT8
precision produce numerically equivalent results: cos\_min within 0.0005
across paths, batch-8 latency within 0.02× when both paths force fp32
fallback. The convergence rules out "wrong tool chain" as an explanation
for failed mixed-precision attempts and pinpoints upstream cumulative
INT8 quantization noise from blocks 0--15 as the binding constraint. To
our knowledge, this is the first explicit cross-tool-chain equivalence
proof in the inference quantization literature.

\textbf{A reproducibility infrastructure pattern for ONNX graph
manipulation testing.} We introduce a project pattern that splits ONNX
graph operations into pure-Python identification and planning helpers
(no \texttt{onnx} dependency) and thin remote-only driver scripts for
actual graph mutation. This pattern enables full unit testing of complex
graph rewriting logic on developer workstations without GPU or native
ONNX dependencies. Our reproducibility kit---271 unit tests across 111
source files, runnable on macOS development hosts---demonstrates the
pattern\textquotesingle s practical value.

\subsection{The remainder of this paper is organized as follows. Section
2 reviews related work in ViT INT8 PTQ, TensorRT mixed-precision
inference, and DPT-style multi-scale fusion. Section 3 describes our
methodology including hardware setup, ONNX export with RoPE source
patching, multi-resolution engine strategy, the 12 precision candidates
evaluated, cosine evaluation protocol, and the pure-Python testing
pattern. Section 4 presents results across the three research questions.
Section 5 discusses the root cause analysis, three-tool-chain
convergence implications, methodological innovations, and comparison to
related work. Section 6 acknowledges limitations including dataset
proxy, single-hardware coverage, and the deferred QAT implementation.
Section 7 concludes. Section 8 summarizes the reproducibility
kit.}\label{the-remainder-of-this-paper-is-organized-as-follows-section-2-reviews-related-work-in-vit-int8-ptq-tensorrt-mixed-precision-inference-and-dpt-style-multi-scale-fusion-section-3-describes-our-methodology-including-hardware-setup-onnx-export-with-rope-source-patching-multi-resolution-engine-strategy-the-12-precision-candidates-evaluated-cosine-evaluation-protocol-and-the-pure-python-testing-pattern-section-4-presents-results-across-the-three-research-questions-section-5-discusses-the-root-cause-analysis-three-tool-chain-convergence-implications-methodological-innovations-and-comparison-to-related-work-section-6-acknowledges-limitations-including-dataset-proxy-single-hardware-coverage-and-the-deferred-qat-implementation-section-7-concludes-section-8-summarizes-the-reproducibility-kit}

\subsection{2. Literature Review}\label{2-literature-review}

We organize prior work into five themes that together establish the
conceptual framework for our investigation: (2.1) the theoretical
underpinnings of post-training quantization (PTQ) and its limits; (2.2)
Vision Transformer (ViT) INT8 PTQ studies and their typical precision
regimes; (2.3) TensorRT mixed-precision inference and the explicit Q/DQ
workflow; (2.4) DPT-style multi-scale fusion and intermediate-layer hook
strategies; and (2.5) a synthesis identifying the specific gap our work
addresses.

\subsubsection{2.1 Theoretical Framework: PTQ vs
QAT}\label{21-theoretical-framework-ptq-vs-qat}

Quantization theory distinguishes between \textbf{post-training
quantization (PTQ)} --- applied to a frozen pretrained model using only
a small calibration dataset --- and \textbf{quantization-aware training
(QAT)} --- which interleaves quantization simulation into the training
loop, allowing model weights to adapt to the quantization grid
{[}Krishnamoorthi, 2018; Nagel et al., 2021{]}. PTQ is attractive for
inference deployment because it requires no retraining and minimal
calibration data, but its precision is fundamentally bounded by the
post-hoc calibration\textquotesingle s ability to fit a fixed weight
distribution into a discrete grid. QAT trades training cost for the
ability to optimize weights jointly with the quantization scheme,
typically recovering 1-3\% additional accuracy on classification
benchmarks at the cost of multi-day training {[}Nagel et al., 2021{]}.

For inference quality assessment of foundation models, \textbf{cosine
similarity against the FP32 baseline} has emerged as a standard fidelity
proxy {[}Touvron et al., 2022; Oquab et al., 2024{]}. Cosine captures
the angular agreement between feature vectors, which is the property
downstream task heads consume. For dense prediction tasks built on
intermediate features (depth estimation, segmentation), the deepest
hooked layer\textquotesingle s cosine fidelity has the most direct
impact on downstream output quality.

\subsubsection{2.2 ViT INT8 PTQ Studies}\label{22-vit-int8-ptq-studies}

Vision Transformer INT8 PTQ has received concentrated attention since
2022. \textbf{GPTQ} {[}Frantar et al., 2023{]} uses approximate
second-order information to optimize per-channel quantization scales,
achieving cos ≥ 0.95 on ViT-S and ViT-B with sub-1\% top-1 accuracy
degradation. \textbf{AWQ} generalizes GPTQ\textquotesingle s insight to
activation-aware weight quantization. \textbf{SmoothQuant} {[}Xiao et
al., 2023{]} addresses the orthogonal problem of activation quantization
difficulty in transformers by transferring scale from activations to
weights via a smoothing parameter α: large activations become
quantization-friendly while weights absorb the corresponding inverse
scale.

Critically, the published ViT INT8 PTQ literature predominantly targets
\textbf{cosine ≥ 0.95 thresholds} suitable for image classification,
where small angular deviations on the final classification logit space
rarely affect top-1 accuracy. For dense prediction downstream consumers
--- DPT {[}Ranftl et al., 2021{]}, MIM-Depth, and follow-up works using
DINOv2/DINOv3 backbones --- empirical observations suggest that
intermediate-feature cos ≥ 0.99 is required to avoid noticeable
degradation in depth estimation and segmentation outputs, but this
stringent regime is rarely benchmarked in the PTQ literature itself. Our
work probes this regime explicitly.

A second observation from the ViT INT8 PTQ literature is that
\textbf{ViT-Large is harder to quantize than ViT-Small or ViT-Base}: the
deeper 24-block stack accumulates more per-block PTQ noise, and the
larger 1024-dimensional hidden state has tighter activation outliers.
Default PTQ on ViT-L/16 frequently collapses to cos ≈ 0.20 {[}our §4.3
observation, consistent with prior anecdotal reports in the SmoothQuant
repository{]}, whereas the same PTQ recipe on ViT-S/B may produce cos ≈
0.97. SmoothQuant partially addresses this by smoothing activation
outliers but does not, by itself, recover cos ≥ 0.99 on
ViT-L/16\textquotesingle s deepest hook.

\subsubsection{2.3 TensorRT and Mixed-Precision
Inference}\label{23-tensorrt-and-mixed-precision-inference}

NVIDIA TensorRT is the de facto inference compiler for production
deployment of vision models on NVIDIA hardware. The 10.x release line
introduced \textbf{explicit Q/DQ ONNX} as the recommended path for INT8
inference, replacing the legacy implicit calibration path that became
deprecated in TRT 10.1 {[}NVIDIA Corporation, 2026{]}. Explicit Q/DQ
embeds quantization scales as ONNX nodes, allowing
TensorRT\textquotesingle s Model Optimizer toolkit to control per-tensor
and per-channel scales precisely.

When PTQ alone is insufficient, TensorRT supports
\textbf{mixed-precision inference} through several mechanisms: (a) the
\texttt{setPrecision()} API allows setting per-layer precision
constraints programmatically; (b) the trtexec command-line
\texttt{-\/-layerPrecisions} flag accepts a name-precision mapping; (c)
the \texttt{-\/-precisionConstraints=obey\textbar{}prefer} flag controls
whether TensorRT must obey or merely prefer the requested precision. The
Model Optimizer\textquotesingle s \texttt{disable\_quantizer} API
provides a fourth, PyTorch-side mechanism to skip Q/DQ insertion for
selected modules during calibration.

The TensorRT documentation describes each of these mechanisms but does
not, to our knowledge, characterize their cross-equivalence under the
explicit Q/DQ workflow. The prior literature has not reported whether
choosing one mechanism over another materially affects mixed-precision
precision-speedup trade-offs, motivating our cross-tool-chain comparison
in §4.4.

A second TensorRT-specific consideration is the \textbf{Myelin pattern
matcher}, the fused-kernel optimization layer in TensorRT 10.x that
handles transformer-specific kernel fusion. Myelin\textquotesingle s
BF16 support has evolved across TRT 10.0--10.13 releases {[}NVIDIA TRT
release notes{]}, and we document a previously unreported BF16 +
explicit Q/DQ incompatibility on the RoPE Constant node in §6.3 specific
to TRT 10.13.

\subsubsection{2.4 DPT-Style Multi-Scale
Fusion}\label{24-dpt-style-multi-scale-fusion}

Vision Transformers\textquotesingle{} single-stride architecture (no
spatial downsampling within the backbone) presents a challenge for dense
prediction tasks that traditionally consume multi-scale features.
\textbf{DPT} {[}Ranftl et al., 2021{]} addresses this by hooking
multiple intermediate transformer blocks and fusing their features via a
learnable fusion head, enabling depth estimation and segmentation
downstream tasks on ViT backbones. For ViT-L/24-block backbones, DPT
recommends hooking blocks \texttt{{[}5,\ 11,\ 17,\ 23{]}} (1-based) ---
equally spaced across the 24-block stack to maximize feature diversity.

Subsequent works using DINOv2 and DINOv3 backbones have largely adopted
DPT\textquotesingle s hook layout or close variants. The empirical
question of whether the equally-spaced DPT layout is optimal --- versus
alternatives that prioritize magnitude balance or earlier-layer
information preservation --- has not been systematically benchmarked on
ViT-L/16 with cos ≥ 0.99 stringent constraint. Our 4-layer ablation in
§4.5 fills this gap.

\subsubsection{2.5 Synthesis and Research
Gap}\label{25-synthesis-and-research-gap}

The five literature themes above converge on a specific gap: while ViT
INT8 PTQ has been studied extensively at cos ≥ 0.95, while TensorRT
mixed-precision has been documented at the API level, and while
DPT-style fusion has been validated at the architectural level, the
\textbf{intersection} of (a) ViT-L/16 + (b) Blackwell sm\_120 + TensorRT
10.13 + (c) cos ≥ 0.99 stringent constraint + (d) cross-tool-chain
mixed-precision comparison has not been mapped. This intersection is
consequential because:

\begin{itemize}
\tightlist
\item
  ViT-L is the deployment-target size for dense prediction tasks
  (smaller backbones underperform on depth/segmentation downstream
  metrics).
\item
  Blackwell sm\_120 (released 2025) introduces 5th-generation Tensor
  Cores with FP8 hardware support that older PTQ studies could not
  benchmark.
\item
  TRT 10.13 (released early 2026) consolidates the explicit Q/DQ
  workflow into the standard PTQ pipeline.
\item
  The cos ≥ 0.99 constraint is empirically required by dense prediction
  downstream consumers and is rarely benchmarked in PTQ literature.
\item
  Cross-tool-chain comparison is critical for engineering decisions but
  absent from prior single-tool-chain studies.
\end{itemize}

\subsection{\texorpdfstring{Our \textbf{conceptual framework} is the
(cosine, speedup) trade-off plane with a "G2 ideal region" defined by
simultaneously meeting cos ≥ 0.99 and speedup ≥ 2.2× thresholds. Within
this framework, we evaluate 12 inference precision candidates spanning
the standard precisions, INT8 partial quantization variants, SmoothQuant
α-sweep, and three independent mixed-precision strategies. We use cosine
convergence across implementations as a falsifiable test: if three
independent tool chains produce equivalent results when forcing the same
blocks back to higher precision, the binding constraint cannot be
tool-chain choice; it must lie elsewhere in the inference path. Section
3 details the methodology, and §4-5 present results within this
framework.}{Our conceptual framework is the (cosine, speedup) trade-off plane with a "G2 ideal region" defined by simultaneously meeting cos ≥ 0.99 and speedup ≥ 2.2× thresholds. Within this framework, we evaluate 12 inference precision candidates spanning the standard precisions, INT8 partial quantization variants, SmoothQuant α-sweep, and three independent mixed-precision strategies. We use cosine convergence across implementations as a falsifiable test: if three independent tool chains produce equivalent results when forcing the same blocks back to higher precision, the binding constraint cannot be tool-chain choice; it must lie elsewhere in the inference path. Section 3 details the methodology, and §4-5 present results within this framework.}}\label{our-conceptual-framework-is-the-cosine-speedup-trade-off-plane-with-a-g2-ideal-region-defined-by-simultaneously-meeting-cos--099-and-speedup--22-thresholds-within-this-framework-we-evaluate-12-inference-precision-candidates-spanning-the-standard-precisions-int8-partial-quantization-variants-smoothquant-ux3b1-sweep-and-three-independent-mixed-precision-strategies-we-use-cosine-convergence-across-implementations-as-a-falsifiable-test-if-three-independent-tool-chains-produce-equivalent-results-when-forcing-the-same-blocks-back-to-higher-precision-the-binding-constraint-cannot-be-tool-chain-choice-it-must-lie-elsewhere-in-the-inference-path-section-3-details-the-methodology-and-4-5-present-results-within-this-framework}

\subsection{3. Methodology}\label{3-methodology}

\subsubsection{3.1 Hardware Setup and Measurement
Protocol}\label{31-hardware-setup-and-measurement-protocol}

All measurements were taken on a single workstation with an NVIDIA
GeForce RTX 5080 GPU (Blackwell architecture, sm\_120 compute
capability, 16 GB VRAM, 300 W TDP), Intel Core Ultra 9 285K CPU, and 128
GB system memory, running Windows 10 Professional. The GPU clock was
locked at 2752 MHz throughout benchmarking via
\texttt{nvidia-smi\ -\/-lock-gpu-clocks=2752}, and
\texttt{trtexec\ -\/-useSpinWait} was used to suppress the Windows
Display Driver Model (WDDM) scheduler\textquotesingle s known 100 ms
timing jitter on consumer-grade Windows GPU stacks. The TensorRT version
under test was 10.13.2.6, paired with CUDA 12.8 and cuDNN 9.x. PyTorch
2.12.0.dev (cu128 nightly) and Transformers 4.57.6 were used for ONNX
export and PTQ calibration.

For latency reporting, we report both trtexec GPU compute time median
and C++ runtime end-to-end median. The trtexec metric isolates pure GPU
kernel execution time excluding host-device transfers, while the C++
runtime metric includes host-to-device input copy, enqueue overhead,
device-to-host output copy, and stream synchronization---reflecting
production deployment cost. We sample 50 measurement iterations after 10
warm-up iterations per (resolution, batch, candidate) configuration.
Trimmed median (middle 50\%) is reported alongside raw median to provide
a jitter-robust comparison.

\subsubsection{3.2 Model and Output
Contract}\label{32-model-and-output-contract}

The model under acceleration is DINOv3 ViT-L/16 LVD-1689M
(\texttt{facebook/dinov3-vitl16-pretrain-lvd1689m}), a 24-block Vision
Transformer with 1024-dimensional hidden states, 16×16 patch size, and
four register tokens {[}Meta AI, 2024{]}. The Hugging Face
implementation by default trims register tokens from intermediate-layer
outputs returned via \texttt{get\_intermediate\_layers()}, yielding a
197-token contract at 224×224 input (1 CLS + 196 patch tokens) which our
project main path adopts for compatibility with downstream DPT-style
fusion heads.

Four intermediate features are exposed as separate ONNX output bindings:
\texttt{feat\_layer\_4}, \texttt{feat\_layer\_12},
\texttt{feat\_layer\_16}, \texttt{feat\_layer\_20} (1-based, equivalent
to 0-based indices \texttt{(3,\ 11,\ 15,\ 19)}). At 224×224 input, each
output has shape \texttt{{[}B,\ 197,\ 1024{]}} per the contract. At
336×336, the patch grid expands to 21×21 yielding 442 tokens; at
518×518, the patch grid is 32×32 yielding 1025 tokens (with
\texttt{floor(518/16)=32}, accepting a 6-pixel input cropping from the
32×16=512 grid).

\subsubsection{3.3 ONNX Export and RoPE
Source-Patch}\label{33-onnx-export-and-rope-source-patch}

ONNX export uses \texttt{torch.onnx.export()} at opset 19. A non-trivial
obstacle is that DINOv3\textquotesingle s RoPE (Rotary Position
Embedding) implementation contains an \texttt{aten::if} conditional
branch in \texttt{angles.tile(2)}, which becomes an ONNX \texttt{If}
node after export. TensorRT 10.13\textquotesingle s
\texttt{IIfConditionalOutputLayer} build path is observed to fail
intermittently on the resulting graph (Issue \#4603/\#4558 in NVIDIA TRT
issue tracker). We resolve this with a source-level patch: replacing
\texttt{angles.tile(2)} with
\texttt{torch.cat((angles,\ angles),\ dim=-1)} in the RoPE forward
function, eliminating the conditional branch entirely while preserving
the mathematical equivalence. This ADR-007 decision is stable across the
entire 24-block stack and produces an \texttt{If}-free ONNX that
TensorRT 10.13 builds without retry.

\subsubsection{3.4 Multi-Resolution Engine
Strategy}\label{34-multi-resolution-engine-strategy}

Per ADR-002 and ADR-009, we use a static-spatial / dynamic-batch profile
strategy: each resolution gets an independent ONNX file and an
independent TensorRT engine binary, with batch as the only dynamic
dimension. Three resolutions (224, 336, 518) are deployed. For the
518×518 batch-8 configuration, the 16 GB VRAM constraint required a
profile of \texttt{min=1,\ opt=4,\ max=8} rather than the standard
\texttt{max=32} used at smaller resolutions; an additional independent
engine with \texttt{max=8} was built to serve the batch-8 measurement
specifically, avoiding profile widening that degrades kernel selection
quality.

Each engine binary is paired with a per-engine TensorRT timing cache
file, persisted to disk to enable build-time consistency across re-runs.
The timing cache is an explicit per-engine artifact rather than a shared
cache, because the optimal kernel choices vary substantially across
precision modes and we observed cross-pollination effects when sharing
caches between engines.

\subsubsection{3.5 Precision Candidates
Evaluated}\label{35-precision-candidates-evaluated}

Twelve precision candidates were evaluated, organized into four
categories.

\emph{Standard precisions (4 candidates):} FP32 (baseline), BF16 prefer
(\texttt{-\/-bf16} with TF32 disabled), FP16 (\texttt{-\/-fp16} with
TF32 disabled), and FP8 default ModelOpt PTQ.

\emph{Node-level partial INT8 (3 candidates):} INT8 ModelOpt explicit
Q/DQ with allow-lists restricted to (a) MatMul layers 16--19, (b) MatMul
layer 19 only, (c) MatMul layer-19-attention only.

\emph{INT8 SmoothQuant α-sweep (3 candidates):} Full-network INT8 Q/DQ
with the SmoothQuant activation-smoothing transformation {[}Xiao et al.,
2023{]} applied at three α settings (0.5, 0.7, 0.8), calibrated on 500
Imagenette images.

\emph{Mixed-precision strategies (3 candidates, all targeting blocks
16--19 → non-INT8):}

\begin{enumerate}
\def\labelenumi{\arabic{enumi}.}
\tightlist
\item
  \textbf{PyTorch ModelOpt \texttt{disable\_quantizer}}: SmoothQuant
  α=0.8 with \texttt{*model.layer.\{16,17,18,19\}.*} wildcards passed to
  the ModelOpt config to skip quantizer insertion.
\item
  \textbf{TensorRT \texttt{-\/-layerPrecisions}}: Full SmoothQuant Q/DQ
  ONNX, with trtexec
  \texttt{-\/-int8\ -\/-precisionConstraints=obey\ -\/-layerPrecisions=\textless{}100\ nodes\textgreater{}:fp32}
  forcing block 16--19 compute-heavy nodes
  (MatMul/Add/LayerNormalization/Softmax) to FP32 at TensorRT
  command-line level.
\item
  \textbf{ONNX-level Q/DQ stripping (V1.2)}: Direct ONNX graph rewriting
  that removes 96 internal Q/DQ nodes (48 pairs × 2) within blocks
  16--19 and rewires the surrounding edges, then trtexec
  \texttt{-\/-int8} build allowing TensorRT to fall back to FP32 for the
  un-quantized region. Per ADR-010 § 4.3, no boundary Q/DQ pairs exist
  between blocks 15 and 16, or between blocks 19 and 20, in SmoothQuant
  α=0.8 ONNX, so a simple delete-and-rewire suffices without boundary
  preservation.
\end{enumerate}

For each candidate, we report (resolution, batch) speedup against FP32
across \texttt{(224,\ \{1,\ 8,\ 32\})}, \texttt{(336,\ \{1,\ 4,\ 8\})},
and \texttt{(518,\ \{1,\ 2,\ 4,\ 8\})}.

\subsubsection{3.6 Cosine Evaluation
Protocol}\label{36-cosine-evaluation-protocol}

For each candidate\textquotesingle s engine, we measure
feat\_layer\_\{4, 12, 16, 20\} cosine similarity against the FP32
baseline engine on 1000 real images at 224×224 input (additional
resolutions for the BF16 prefer candidate). We compute per-image cosine
for each of the four outputs, then aggregate to cos\_mean (arithmetic
mean across the 1000 images) and cos\_min (minimum across the 1000
images). We report cos\_min as the primary precision metric because
production deployment is sensitive to worst-case feature drift;
cos\_mean provides aggregate fidelity context.

The image source is Imagenette2-320 validation split {[}fast.ai{]}, of
which 1000 images are used for cosine evaluation and 500 disjoint images
are used for SmoothQuant calibration---the two sets are explicitly
stratified-sampled and mutually exclusive to prevent calibration
leakage. Imagenette serves as a proxy for ImageNet val 50K, which is
gated on Hugging Face under \texttt{ILSVRC/imagenet-1k} and was
inaccessible during the project window (\texttt{403\ GatedRepoError}).
While Imagenette is a strict 10-class subset of
ImageNet\textquotesingle s 1000 classes, the smaller class diversity may
underestimate worst-case PTQ error; we acknowledge this in § 6 and
provide a one-command swap-in path
(\texttt{scripts/export\_hf\_imagenet\_parquet\_images.py}) for ImageNet
val once authorization is obtained.

\subsubsection{3.7 Cross-Language Parity
Methodology}\label{37-cross-language-parity-methodology}

To verify that the C++ runtime wrapper used in production deployment
produces identical numerical outputs to the Python TensorRT runtime used
during research, we feed both with the same deterministic sine input (a
closed-form pixel-value tensor reproducible from a single random seed)
and compare per-output numerical agreement. The comparison metrics are
max absolute error, RMSE, and cosine similarity---the cosine of an
output against itself across the two runtimes.

The C++ runtime is a thin RAII wrapper around the TensorRT 10.13.2.6 C++
API, built with Microsoft Visual C++ (MSVC) 19.44 under Visual Studio
2022\textquotesingle s Developer Command Prompt. We discovered during
project development that MinGW g++ 14.2 produces an ABI-incompatible
binary---\texttt{getIOTensorName()} returns garbage---forcing exclusive
use of MSVC for the CUDA-coupled C++ binary. The CMake + Ninja build
system produces a single executable that loads any project ONNX engine
and runs inference; the same executable, the same engine binary, and the
same deterministic input are then fed into a Python TRT runtime. We
require max\_abs\_error = 0, RMSE = 0, and cosine = 1.0 across all four
outputs---not epsilon-close, but bit-identical.

This protocol is run for every (resolution, candidate) pair in the
project\textquotesingle s main path, including FP32, BF16 prefer,
partial INT8, layer-19 INT8 at 224×224, and FP32 / BF16 prefer at
336×336 and 518×518. All twelve resulting parity reports show
bit-identical outputs.

\subsubsection{3.8 Pure-Python Testing Pattern for ONNX Graph
Manipulation}\label{38-pure-python-testing-pattern-for-onnx-graph-manipulation}

A methodological contribution arising from the V1.2 mixed-precision
implementation (§ 3.5 third strategy) is a project pattern that splits
ONNX graph manipulation into two independent layers:

\begin{enumerate}
\def\labelenumi{\arabic{enumi}.}
\item
  \textbf{Pure-Python identification and planning helpers} (no
  \texttt{onnx} library dependency). Operating on simple
  \texttt{(name,\ op\_type,\ inputs,\ outputs)} tuples projected from
  \texttt{onnx.GraphProto.node}, these helpers identify Q/DQ pairs by
  tensor edge connectivity, classify each pair as
  internal/boundary\_input/boundary\_output relative to a requested
  block range, and emit a \texttt{StripPlan} data structure describing
  exactly which nodes to delete and which downstream tensor references
  to rewire. The planner is fully unit-testable on synthetic mini
  graphs.
\item
  \textbf{Thin remote-only driver scripts} that import \texttt{onnx},
  project the graph, invoke the pure-Python planner, and apply the
  resulting plan via direct \texttt{onnx.GraphProto} mutation followed
  by \texttt{onnx.save()}.
\end{enumerate}

This split enables developer-side unit testing of complex graph
rewriting logic without GPU or native ONNX dependencies, in our case
running 271 unit tests across 111 source files on a macOS development
host while the actual graph mutation runs on the Windows + RTX 5080
deployment node. The pattern has produced unit-tested implementations of
three V1.1/V1.2 modules: \texttt{layer\_precision.py} (per-layer
precision argument generation, round 17),
\texttt{onnx\_qdq\_stripper.py} (Q/DQ pair identification and
classification, round 24), and \texttt{onnx\_qdq\_strip\_planner.py}
(Q/DQ strip planning with conflict detection, round 25).

\subsubsection{3.9 Reproducibility
Infrastructure}\label{39-reproducibility-infrastructure}

To support community reproducibility, we publish a complete artifact
kit. The benchmark matrix is a 56-row machine-readable CSV with one row
per (runtime, candidate, reference\_precision, batch\_size, resolution)
tuple, enumerating trtexec GPU compute time speedup and throughput
speedup. Eight SVG figures cover speedup bar charts, multi-resolution
cosine plots, the 12-point INT8 sensitivity scatter, and the 4-layer
ablation diversity-versus-balance plot. A unified figure regeneration
entry point (\texttt{scripts/build\_all\_figures.py\ -\/-allow-missing})
regenerates all 8 SVGs from the underlying data and emits a top-level
\texttt{figures\_index.json} cross-subsystem index for diff-friendly
verification across re-runs.

A SHA256 manifest covering 419+ artifact files (ONNX, engines, timing
caches, JSON reports, SVG figures, log files, etc.) is generated by
\texttt{check\_assets.py\ -\/-output\ PATH}, which performs an atomic
write (tempfile + \texttt{os.replace}) and includes the target file path
in its exclude set. This eliminates a subtle pre-existing bug where
shell-redirected manifest generation produced 0-byte target files that
the manifest scanner then recorded as stale entries---the so-called
"manifest-generates-itself" 0-byte file problem.

\subsection{\texorpdfstring{The \texttt{sync\_remote\_windows\_repo.py}
script provides bidirectional macOS ↔ Windows synchronization; the
default direction pushes the source tree to the Windows GPU host, while
\texttt{-\/-pull-reports} performs reverse-pull of text-only artifacts
(\texttt{.json}, \texttt{.md}, \texttt{.csv}, \texttt{.svg},
\texttt{.png}, \texttt{.jpg}, \texttt{.log}, \texttt{.txt}) via
PowerShell-packed zip archives, sidestepping flaky scp behavior on the
cpolar SSH tunnel. This separation between large binary artifacts (kept
on the GPU host) and text-only reproducible artifacts (round-tripped to
the development host) is a project-wide design choice that simplifies
day-to-day workflows without compromising
reproducibility.}{The sync\_remote\_windows\_repo.py script provides bidirectional macOS ↔ Windows synchronization; the default direction pushes the source tree to the Windows GPU host, while -\/-pull-reports performs reverse-pull of text-only artifacts (.json, .md, .csv, .svg, .png, .jpg, .log, .txt) via PowerShell-packed zip archives, sidestepping flaky scp behavior on the cpolar SSH tunnel. This separation between large binary artifacts (kept on the GPU host) and text-only reproducible artifacts (round-tripped to the development host) is a project-wide design choice that simplifies day-to-day workflows without compromising reproducibility.}}\label{the-sync_remote_windows_repopy-script-provides-bidirectional-macos--windows-synchronization-the-default-direction-pushes-the-source-tree-to-the-windows-gpu-host-while---pull-reports-performs-reverse-pull-of-text-only-artifacts-json-md-csv-svg-png-jpg-log-txt-via-powershell-packed-zip-archives-sidestepping-flaky-scp-behavior-on-the-cpolar-ssh-tunnel-this-separation-between-large-binary-artifacts-kept-on-the-gpu-host-and-text-only-reproducible-artifacts-round-tripped-to-the-development-host-is-a-project-wide-design-choice-that-simplifies-day-to-day-workflows-without-compromising-reproducibility}

\subsection{4. Results}\label{4-results}

We present results in five subsections corresponding to the three
research questions and the methodological cross-validation. Section 4.1
establishes the BF16 prefer baseline with multi-resolution speedup
measurements (RQ1). Section 4.2 reports the BF16 prefer cosine fidelity
across resolutions (RQ1). Section 4.3 displays the comprehensive
12-candidate sensitivity scatter mapping the operating region (RQ1).
Section 4.4 presents the three-tool-chain mixed-precision convergence
proof (RQ2 + RQ3). Section 4.5 reports the 4-layer ablation
diversity-vs-balance trade-off. Section 4.6 verifies cross-language
Python ↔ C++ runtime parity as a sanity check.

\subsubsection{4.1 BF16 Prefer Speedup
(RQ1)}\label{41-bf16-prefer-speedup-rq1}

Table 1 reports trtexec GPU compute time median speedup of BF16 prefer
engines against the FP32 baseline at all locked-clock + spin-wait
measurement configurations.

\textbf{Table 1.} BF16 prefer trtexec GPU compute median speedup vs FP32
baseline. Locked GPU clock 2752 MHz, \texttt{-\/-useSpinWait} enabled,
50 iterations after 10-iteration warm-up.

{\def\LTcaptype{none} % do not increment counter
\begin{longtable}[]{@{}lrrrrr@{}}
\toprule\noalign{}
Resolution & b1 & b2 & b4 & b8 & b32 \\
\midrule\noalign{}
\endhead
\bottomrule\noalign{}
\endlastfoot
224 × 224 & 2.45× & --- & --- & 2.81× & 3.25× \\
336 × 336 & 2.80× & --- & 2.96× & 3.25× & --- \\
518 × 518 & 3.12× & 3.50× & 3.76× & \textbf{3.86×} & --- (16 GB cap) \\
\end{longtable}
}

The peak observation is \textbf{3.86× speedup at r518 batch 8}. The 16
GB VRAM constraint prevents larger batch sizes at r518; r336 is capped
at b8 by the project\textquotesingle s static spatial profile design
(extending to b16/b32 would require independent engines). Within each
resolution, speedup increases monotonically with batch size as expected
for compute-bound BF16 kernels with reduced amortized overhead per
inference. Across resolutions at fixed batch, larger inputs achieve
higher speedup because the FP32 baseline\textquotesingle s compute cost
grows superlinearly with token count (the BF16 kernel maintains better
arithmetic intensity at high token counts on Blackwell sm\_120 5th-gen
Tensor Cores).

Figure 1 (\texttt{figures/benchmark\_trtexec\_bf16\_speedup.svg})
visualizes the multi-resolution trtexec speedup as a grouped bar chart
with separate panels per resolution.

\textbf{Table 2} reports C++ runtime end-to-end median speedup, which
includes host-to-device input copy, enqueue overhead, device-to-host
output copy, and stream synchronization---reflecting production
deployment cost.

\textbf{Table 2.} C++ runtime end-to-end median speedup vs FP32 baseline
(50 iter × 10 warm-up).

{\def\LTcaptype{none} % do not increment counter
\begin{longtable}[]{@{}lrr@{}}
\toprule\noalign{}
Resolution & b1 & b8 \\
\midrule\noalign{}
\endhead
\bottomrule\noalign{}
\endlastfoot
224 × 224 & 2.27× & 2.50× \\
336 × 336 & 2.60× & 2.85× \\
518 × 518 & 2.83× & \textbf{3.40×} \\
\end{longtable}
}

C++ end-to-end speedup is consistently 0.18--0.46× lower than trtexec
GPU-compute-only speedup, attributable to fixed host-side overhead
(copies + enqueue) that does not scale with the precision change. The
peak C++ speedup (3.40× at r518 b8) remains substantially above the G2 ≥
2.2× threshold and represents the project\textquotesingle s headline
production-grade result. Figure 2
(\texttt{figures/benchmark\_cpp\_runtime\_speedup.svg}) plots these
results.

\subsubsection{4.2 BF16 Prefer Cosine Fidelity
(RQ1)}\label{42-bf16-prefer-cosine-fidelity-rq1}

Table 3 reports per-output cos\_min and cos\_mean against the FP32
baseline on 1000 Imagenette images at all three resolutions.

\textbf{Table 3.} BF16 prefer feat\_layer\_\{4, 12, 16, 20\} cosine
similarity vs FP32 baseline (Imagenette 1000-image evaluation, batch 8).

{\def\LTcaptype{none} % do not increment counter
\begin{longtable}[]{@{}lrrrr@{}}
\toprule\noalign{}
Resolution & feat\_layer\_4 cos\_min & feat\_layer\_12 cos\_min &
feat\_layer\_16 cos\_min & feat\_layer\_20 cos\_min \\
\midrule\noalign{}
\endhead
\bottomrule\noalign{}
\endlastfoot
224 × 224 & 0.999933 & 0.999664 & 0.998943 & 0.998749 \\
336 × 336 & 0.999891 & 0.999276 & 0.998394 & 0.998493 \\
518 × 518 & 0.999868 & 0.999075 & 0.998604 & \textbf{0.999171} \\
\end{longtable}
}

All twelve cells satisfy cos\_min ≥ 0.998, comfortably above the G2
cosine threshold of 0.99. A counter-intuitive finding emerges at the
deepest hooked layer: r518 feat\_layer\_20 cos\_min (0.999171) exceeds
r224 feat\_layer\_20 cos\_min (0.998749). This contradicts the naive
expectation that deeper layers and longer token sequences amplify
quantization error. We attribute the inversion to \textbf{patch-token
dilution}: at r518, feat\_layer\_20 has 1024 patch tokens (plus 1 CLS)
versus r224\textquotesingle s 196, so the per-token BF16 quantization
error is averaged over a 5.2× larger token population, reducing the
worst-case sample-level cosine deviation. This is a
hardware-dataset-specific observation we have not seen reported in prior
literature. Figures 3 and 4
(\texttt{figures/benchmark\_bf16\_cosine\_min.svg} and
\texttt{benchmark\_bf16\_cosine\_mean.svg}) visualize per-output cosine
across resolutions.

\subsubsection{4.3 12-Candidate Sensitivity Map (RQ1, central
result)}\label{43-12-candidate-sensitivity-map-rq1-central-result}

Figure 5 (\texttt{figures/benchmark\_bf16\_vs\_int8\_tradeoff.svg}) is
the project\textquotesingle s central visualization: a 12-point scatter
with X-axis = feat\_layer\_20 cos\_mean, Y-axis = trtexec batch-8
latency speedup vs FP32 baseline, and the G2 ideal region (cos ≥ 0.99 ∧
speedup ≥ 2.2×) shaded in green. Each point is color-coded by candidate
family (BF16 in blue, FP8 default in orange, INT8 partial in
pink/yellow, SmoothQuant α-sweep in olive, mixed-precision strategies in
green/gray).

\textbf{Key observations from the scatter:}

\begin{enumerate}
\def\labelenumi{\arabic{enumi}.}
\tightlist
\item
  \textbf{BF16 prefer is the only candidate inside the G2 ideal region.}
  It occupies the upper-right corner with cos\_mean ≈ 0.9998 and speedup
  ≈ 2.81× (b8 r224 datapoint).
\item
  \textbf{FP8 default ModelOpt PTQ achieves the highest raw speedup}
  (5.05× vs FP32 at b32, 1.55× vs BF16 prefer at b32) but cos\_mean
  collapses to 0.20 at feat\_layer\_20 --- symmetric failure to default
  INT8.
\item
  \textbf{The trade-off curve "the smaller the quantization range, the
  higher the cosine, the lower the speedup"} is visually unambiguous
  along the INT8 partial sequence: layer-19-attention-only (cos 0.99941
  / speed 1.04×), layer 19 only (cos 0.99566 / speed 1.07×), layers
  16-19 (cos 0.989 / speed 1.22×).
\item
  \textbf{SmoothQuant α=0.8 best} (cos 0.982 / speed 3.48×) is the
  closest to the G2 region of all INT8 candidates, but its cos\_min of
  0.968 falls 0.022 short of the 0.99 cosine threshold.
\item
  \textbf{Three mixed-precision strategies (PyTorch / TRT / ONNX)}
  cluster tightly around (cos 0.984, speed 2.4×), all outside the G2
  region --- discussed in detail in §4.4.
\end{enumerate}

The scatter directly answers RQ1: among 12 standard inference precision
candidates, only BF16 prefer satisfies (cos ≥ 0.99 ∧ speedup ≥ 2.2×) on
this model + hardware combination at the current TRT 10.13 release.

\subsubsection{4.4 Three-Tool-Chain Mixed-Precision Convergence (RQ2 +
RQ3)}\label{44-three-tool-chain-mixed-precision-convergence-rq2--rq3}

Table 4 reports the three independent mixed-precision implementations
targeting blocks 16--19 → non-INT8.

\textbf{Table 4.} Three-tool-chain mixed-precision results (blocks
16--19 → non-INT8). All build on the same SmoothQuant α=0.8 calibration;
only the mixed-precision implementation layer differs.

{\def\LTcaptype{none} % do not increment counter
\begin{longtable}[]{@{}llrrr@{}}
\toprule\noalign{}
Tool chain & Implementation layer & feat\_layer\_20 cos\_min &
feat\_layer\_20 cos\_mean & b8 trtexec speedup vs FP32 \\
\midrule\noalign{}
\endhead
\bottomrule\noalign{}
\endlastfoot
Full SmoothQuant α=0.8 (no mixed-precision) & (baseline) & 0.968 & 0.982
& 3.48× \\
ModelOpt \texttt{disable\_quantizer} skip 16-19 & PyTorch (calibration
time) & 0.971 (+0.003) & 0.984 (+0.002) & 2.41× (-31\%) \\
trtexec \texttt{-\/-layerPrecisions=l16-19:fp32} & TRT (build time,
command-line) & 0.9683 (≈) & 0.9822 (≈) & 3.43× (-1.4\%) \\
\textbf{ONNX-level Q/DQ stripping (V1.2)} & \textbf{ONNX library (graph
rewrite)} & \textbf{0.9705 (+0.003)} & \textbf{0.9842 (+0.002)} &
\textbf{2.39× (-31\%)} \\
\end{longtable}
}

The convergence is striking: across three independent implementations
operating at three distinct levels of the inference stack (PyTorch
calibration, TRT command-line, ONNX library), the cos\_min values span
only \textbf{0.0027} (0.9683 to 0.971) and the b8 speedup values span
\textbf{0.04×} (2.39× to 2.43×) when the implementation actually forces
FP32 fallback for blocks 16--19. The TRT \texttt{-\/-layerPrecisions}
row shows nearly identical values to the SmoothQuant baseline ---
TensorRT does not honor the precision constraint when explicit Q/DQ is
present in the ONNX, effectively making this strategy a no-op as
discussed in §5.

The two strategies that successfully force FP32 in blocks 16--19
(ModelOpt disable + ONNX strip) yield equivalent results: cos\_min 0.971
vs 0.9705 (within 0.0005), b8 speedup 2.41× vs 2.39× (within 0.02×).
This convergence proof eliminates "wrong tool chain" as a hypothesis for
the failure to enter the G2 ideal region --- the binding constraint must
lie elsewhere in the inference path (RQ3).

\subsubsection{4.5 4-Layer Hook Selection
Ablation}\label{45-4-layer-hook-selection-ablation}

Table 5 reports the empirical comparison of three 4-layer hook selection
strategies on the same DINOv3 ViT-L/16 backbone.

\textbf{Table 5.} 4-layer hook selection ablation (1000 Imagenette
images, PyTorch hooks).

{\def\LTcaptype{none} % do not increment counter
\begin{longtable}[]{@{}llrr@{}}
\toprule\noalign{}
Candidate & Layers (1-based) & Mean inter-output cosine & Max/min
magnitude ratio \\
\midrule\noalign{}
\endhead
\bottomrule\noalign{}
\endlastfoot
project (selected) & 4 / 12 / 16 / 20 & 0.383 & \textbf{12.6×} (most
balanced) \\
dpt (paper recommended) & 5 / 11 / 17 / 23 & \textbf{0.299} (most
diverse) & 31.9× \\
late-heavy & 6 / 12 / 18 / 24 & 0.339 & 84× (last hook dominates) \\
\end{longtable}
}

The project\textquotesingle s choice (\texttt{{[}4,\ 12,\ 16,\ 20{]}})
achieves the \textbf{tightest magnitude balance} (max/min ratio 12.6×)
at the cost of slightly higher inter-output cosine compared to the DPT
paper\textquotesingle s recommendation (\texttt{{[}5,\ 11,\ 17,\ 23{]}},
mean cosine 0.299 most diverse). The late-heavy variant
(\texttt{{[}6,\ 12,\ 18,\ 24{]}}) exhibits an 84× magnitude imbalance
because feat\_layer\_24\textquotesingle s L2 magnitude grows sharply at
the final transformer block, dominating the multi-scale fusion
contribution from earlier layers. Figure 6
(\texttt{figures/layer\_ablation\_diversity\_vs\_balance.svg}) plots all
three candidates on a 2D space (X = mean inter-output cosine, Y = log10
magnitude max/min), color-coded by candidate (project blue, dpt green,
late red).

We interpret the project\textquotesingle s choice as a
\textbf{diversity-magnitude trade-off} rather than a strict diversity
maximizer: although DPT-style fusion benefits from inter-output
diversity, magnitude imbalance imposes practical cost --- multi-scale
fusion heads must learn to suppress dominant features, increasing
trainable parameter count and degrading data efficiency. The
project\textquotesingle s selection sacrifices \textasciitilde22\%
inter-output diversity (0.383 vs 0.299) for a \textasciitilde2.5×
tighter magnitude balance (12.6× vs 31.9×).

\subsubsection{4.6 Cross-Language Parity (Sanity
Check)}\label{46-cross-language-parity-sanity-check}

Twelve cross-language Python ↔ C++ runtime parity reports were generated
covering (224, 336, 518) × (FP32, BF16 prefer) primary candidates plus
three additional (224, partial INT8 / layer-19 INT8) configurations from
V1.1 sensitivity exploration. All 12 reports show
\textbf{max\_abs\_error = 0, RMSE = 0, cosine = 1.0} across all four
feat\_layer outputs against deterministic sine input, confirming that
the C++ runtime wrapper produces bit-identical numerical outputs to the
Python TRT runtime when sharing the same engine binary and the same
input.

This cross-language parity is a stronger guarantee than typical
"epsilon-close" comparisons: it confirms the project\textquotesingle s
Python research stack and the C++ production runtime are not merely
"approximately equivalent" but exactly equivalent, removing C++ runtime
as a potential confounding source for any precision discrepancies
observed in §4.1--4.5.

\subsubsection{4.7 Throughput-Oriented Serving and Pool Saturation
(V1.0.3
Extension)}\label{47-throughput-oriented-serving-and-pool-saturation-v103-extension}

The V1.0.0 study mapped single-stream BF16 prefer latency speedup but
did not characterize the model\textquotesingle s aggregate throughput
ceiling under concurrent load. Production serving systems aggregate
multiple inference requests, both within a single process (multi-context
dispatch) and across processes (request batching at the serving layer).
Understanding where DINOv3 ViT-L/16 saturates the GPU under each regime
is necessary to determine whether further dense-precision optimization
is feasible --- a question we answer empirically in this section.

\paragraph{4.7.1 Multi-Context Pool
Architecture}\label{471-multi-context-pool-architecture}

We implemented an in-process multi-context pool
(\texttt{TRTInfererPool}) that deserializes the TensorRT engine once at
pool construction and instantiates \emph{N} independent
\texttt{IExecutionContext} objects, each bound to its own
\texttt{cudaStream\_t} and device-side input / output buffers.
Permit-gated capacity (\texttt{std::counting\_semaphore}) ensures the
pool never over-subscribes its slots; per-slot mutexes provide defense
in depth against undocumented thread-safety behavior in the underlying
TensorRT runtime. Round-robin atomic dispatch routes incoming inferences
to the next available slot.

A first-cut implementation deserialized the engine \emph{per slot}
(rather than sharing it across contexts) and crashed at concurrent N≥2
with a TensorRT 10.13 error from
\texttt{MyelinRunnerBase::executeMyelinGraph}
(\texttt{runtime/myelin/runner.cpp:778}). We initially hypothesized
fundamental thread-safety failure in TensorRT 10.13, but two focused
experiments contradicted this: (a) two completely independent OS
processes succeed concurrently at 1.087× over baseline (heavily
diminished throughput due to inter-process CUDA context contention but
no Myelin error), and (b) a single shared engine with two independent
execution contexts on independent streams reproduces concurrent N=2
execution successfully at 1.515× over baseline --- matching the speedup
obtained by Python\textquotesingle s \texttt{tensorrt} bindings under
the GIL. This evidence shows the original crash was caused by duplicate
engine deserialization within the same process, \emph{not} by a
fundamental TensorRT thread-safety constraint. NVIDIA\textquotesingle s
documented "single engine, multiple contexts on independent streams =
thread-safe" architecture is therefore empirically correct at TensorRT
10.13 for this model.

\paragraph{4.7.2 Throughput
Measurement}\label{472-throughput-measurement}

We characterized end-to-end aggregate throughput at three resolutions ×
multiple batch sizes × multiple pool capacities. End-to-end measurements
include the full host-to-device transfer, \texttt{enqueueV3} dispatch,
four device-to-host transfers (one per intermediate-feature output), and
the synchronous stream-wait at call exit, matching the production
serving path.

At the small-load regime (r224, b=1) the V1.0.2 single-stream baseline
delivers 343 qps. The multi-context pool at N=2 with batch=8 reaches 644
qps (1.875×); over-saturation by allowing more caller threads than pool
slots (N=4 slots × 16 caller threads, with caller-side pinned host
memory via \texttt{cudaMallocHost}) raises the empirical ceiling to
\textbf{720 qps (2.094× over the V1.0.2 baseline; 5.24× over the
original FP32 trtexec gpu-compute baseline of 137 qps at r224 b=1)}. The
plan target of 800 qps for this regime was set assuming inter-request
batching at a serving layer (e.g., Triton\textquotesingle s
\texttt{dynamic\_batching}); without that layer the in-process pool
plateaus around 720 qps because per-call \texttt{cudaStreamSynchronize}
serializes within each slot regardless of caller concurrency. A
pure-compute prototype (\texttt{test\_concurrent\_contexts.cpp}) that
omits host transfers reaches 922 qps (2.683× over baseline), confirming
that the residual \textasciitilde22\% gap is bounded by host-device
transfer rather than by GPU compute.

At the medium-load regime (r336, b=8), the baseline of 363 qps lifts to
only 401 qps (1.10×) under N=2 multi-context. At the saturation regime
(r518, b=8), the baseline of 156 qps moves to 160 qps (1.02×). Both
reproduce the V1.0.2 prior finding that larger batches and resolutions
leave essentially no GPU compute headroom for additional contexts to
fill. Larger batch sizes (b=16, b=32) regress sharply at r336 and r518
because the model crosses into a memory-bound regime where HBM bandwidth
becomes the binding constraint --- at r518 b=16 throughput collapses to
34 qps with \textasciitilde400 ms per-call latency.

\paragraph{4.7.3 GPU Utilization Saturation
Evidence}\label{473-gpu-utilization-saturation-evidence}

We recorded SM utilization, HBM utilization, board power, and
temperature at 100 ms cadence via \texttt{nvidia-smi\ -\/-query-gpu=...}
during each benchmark run, aggregated to mean / p50 / p95 / max
statistics. Across the four primary regimes (Table V1.0.3-A):

\textbf{Table V1.0.3-A. GPU utilization across BF16-prefer regimes (RTX
5080, locked clock 2752 MHz).}

{\def\LTcaptype{none} % do not increment counter
\begin{longtable}[]{@{}llllll@{}}
\toprule\noalign{}
Regime & mean SM \% & p50 SM \% & p95 SM \% & mean power (W) & peak
power (W) \\
\midrule\noalign{}
\endhead
\bottomrule\noalign{}
\endlastfoot
r518 b8 N=1 (saturation) & 99.08 & 99 & 100 & 326 & 333 \\
r336 b8 N=1 (medium load) & 96.39 & 99 & 99 & 289 & 293 \\
r224 b1 N=1 (low load) & 88.24 & 92 & 92 & 192 & 194 \\
r224 b1 N=2 (multi-stream) & 95.77 & 99 & 99 & 258 & 263 \\
r224 b8 N=4 t=16 (G1 production) & \textbf{98.74} & 99 & \textbf{100} &
292 & 316 \\
\end{longtable}
}

The high-batch / high-resolution regimes are physically saturated at
single-stream --- multi-context cannot extract additional throughput
because there are no idle SMs to fill. Low-batch regimes have
\textasciitilde12 percentage points of SM idle headroom that
multi-context successfully fills (r224 b1: 88.24\% → 95.77\% under N=2).
The G1 production configuration that achieves the 720 qps end-to-end
ceiling drives SM utilization to 98.74\% mean and 100\% p95 --- within
1.3 percentage points of the per-pixel hardware ceiling. Power
consumption stays within 92\% of the 360 W RTX 5080 board limit across
all regimes (peak 332.6 W during sustained r518 b=8); temperature peaks
at 72 °C, ten degrees below the \textasciitilde83 °C thermal-throttle
threshold. Tensor Core utilization across BF16 dense inference paths is
structurally bounded near 70\% by the BF16 path\textquotesingle s design
--- INT8 / FP8 / 2:4-sparse Tensor Core paths require V1.3 QAT to
unlock.

\paragraph{4.7.4 Concurrent-Execution
Bit-Exactness}\label{474-concurrent-execution-bit-exactness}

We verified output equivalence between the multi-context pool and the
single-context reference using a dedicated validator
(\texttt{pool\_cos\_validator.exe}), submitting the same deterministic
input through both paths and comparing each of the four feature-layer
outputs element-wise. At r224 b=1 N=2 (10 iterations × 2 caller threads
= 20 concurrent inferences) and at r224 b=8 N=2 (the G1 production batch
size, 20 concurrent inferences), every output passes a strict bit-exact
check: \texttt{max\_abs\_error\ =\ 0} across all four outputs, cosine
similarity = 1.000 (perfect, far exceeding the 0.997 V1.0.3 acceptance
threshold). The multi-context pool produces deterministic, race-free
output identical to the single-context reference; concurrent execution
introduces no precision drift.

\paragraph{4.7.5 Implications for
Quantization}\label{475-implications-for-quantization}

\subsection{\texorpdfstring{The combined evidence --- full SM saturation
across two of three resolution regimes (96--99\% at r336 b8 and r518 b8
single-stream), a 720 qps in-process throughput ceiling at the third
regime (98.74\% SM at G1 production config), perfect bit-exact
concurrent execution, and (from §4.3) a confirmed PTQ precision wall
across INT8, FP8, and 2:4-sparse paths --- establishes that further
acceleration of this model on this hardware requires \textbf{reducing
per-inference compute}, not extracting more from the existing dense BF16
path. The runtime axis (saturation) and the precision axis (PTQ wall)
jointly identify quantization-aware fine-tuning (QAT) as the only
published technique that can simultaneously hold cosine similarity above
the 0.99 hard gate while replacing dense BF16 matmuls with INT8 / FP8 /
2:4-sparse Tensor Core operations. The V1.0.3 utilization data thus
provides the quantitative empirical motivation for the V1.3 milestone
discussed in
§5.3.}{The combined evidence --- full SM saturation across two of three resolution regimes (96--99\% at r336 b8 and r518 b8 single-stream), a 720 qps in-process throughput ceiling at the third regime (98.74\% SM at G1 production config), perfect bit-exact concurrent execution, and (from §4.3) a confirmed PTQ precision wall across INT8, FP8, and 2:4-sparse paths --- establishes that further acceleration of this model on this hardware requires reducing per-inference compute, not extracting more from the existing dense BF16 path. The runtime axis (saturation) and the precision axis (PTQ wall) jointly identify quantization-aware fine-tuning (QAT) as the only published technique that can simultaneously hold cosine similarity above the 0.99 hard gate while replacing dense BF16 matmuls with INT8 / FP8 / 2:4-sparse Tensor Core operations. The V1.0.3 utilization data thus provides the quantitative empirical motivation for the V1.3 milestone discussed in §5.3.}}\label{the-combined-evidence--full-sm-saturation-across-two-of-three-resolution-regimes-9699-at-r336-b8-and-r518-b8-single-stream-a-720-qps-in-process-throughput-ceiling-at-the-third-regime-9874-sm-at-g1-production-config-perfect-bit-exact-concurrent-execution-and-from-43-a-confirmed-ptq-precision-wall-across-int8-fp8-and-24-sparse-paths--establishes-that-further-acceleration-of-this-model-on-this-hardware-requires-reducing-per-inference-compute-not-extracting-more-from-the-existing-dense-bf16-path-the-runtime-axis-saturation-and-the-precision-axis-ptq-wall-jointly-identify-quantization-aware-fine-tuning-qat-as-the-only-published-technique-that-can-simultaneously-hold-cosine-similarity-above-the-099-hard-gate-while-replacing-dense-bf16-matmuls-with-int8--fp8--24-sparse-tensor-core-operations-the-v103-utilization-data-thus-provides-the-quantitative-empirical-motivation-for-the-v13-milestone-discussed-in-53}

\subsection{5. Discussion}\label{5-discussion}

We discuss four aspects of our findings: (5.1) the root cause of the
empirical operating-region boundary identified in §4.3; (5.2) the
implications of the three-tool-chain convergence proof of §4.4 for
engineering practice; (5.3) what our results --- together with the
V1.0.3 §4.7 throughput-saturation evidence --- suggest for the V1.3
quantization-aware fine-tuning (QAT) future-work path; (5.4) the
methodological pattern we identify for cross-tool-chain validation in
inference systems research; and (5.5) how our findings relate to and
extend prior work in ViT INT8 PTQ and TensorRT mixed-precision
literature.

\subsubsection{5.1 Root Cause: Upstream Cumulative Quantization
Noise}\label{51-root-cause-upstream-cumulative-quantization-noise}

The central empirical finding of this study is that no INT8 PTQ
candidate---including all SmoothQuant α settings and all three
mixed-precision strategies that force blocks 16--19 back to
FP32---enters the G2 ideal region (cos ≥ 0.99 ∧ speedup ≥ 2.2×) on this
model + hardware combination. The convergence of the three independent
mixed-precision tool chains in §4.4 (cos\_min span 0.0027, speedup span
0.04×) refutes the hypothesis that the failure is due to tool-chain
implementation choice. Instead, we attribute the failure to
\textbf{upstream cumulative INT8 quantization noise from blocks 0--15}.

To quantify the mechanism: each transformer block introduces a per-block
INT8 PTQ error that propagates forward. With SmoothQuant α=0.8, the
per-block deviation against the FP32 baseline is empirically
\textasciitilde10⁻²·⁵ in feat\_layer\_4 cos\_min (0.991), accumulating
to \textasciitilde10⁻¹·⁵ in feat\_layer\_20 cos\_min (0.968) --- a
deviation that \textbf{already exceeds the 10⁻² gap required to meet the
0.99 cosine threshold by the time the activation reaches block 16}.
Forcing blocks 16--19 to FP32 internal computation cannot recover
information already lost upstream: the FP32 kernels operate on
contaminated input and produce correspondingly contaminated output. This
is consistent with the V1.1 + V1.2 mixed-precision experimental
observation (Table 4) that all three tool chains converge to
feat\_layer\_20 cos\_min ≈ 0.97, \textasciitilde22\% short of the cosine
threshold.

The mechanism explains why mixed-precision \emph{targeting the wrong
locus} is insufficient: the binding constraint is upstream weight
calibration error, not downstream computation precision. Closing the
cosine gap thus requires intervening at the weight level --- i.e.,
quantization-aware fine-tuning (QAT) rather than further variations on
PTQ.

\subsubsection{5.2 Implications of Three-Tool-Chain
Convergence}\label{52-implications-of-three-tool-chain-convergence}

The empirical convergence of three independent mixed-precision
implementations (PyTorch ModelOpt \texttt{disable\_quantizer}, TensorRT
\texttt{-\/-layerPrecisions}, ONNX-level Q/DQ stripping) is a
methodological contribution beyond the negative result it establishes.
Three engineering observations follow.

\textbf{First, tool-chain choice is irrelevant for explicit Q/DQ
workflows.} When TensorRT receives an explicit Q/DQ ONNX, the build-time
Q/DQ nodes constrain the kernel selection regardless of the upstream
toolchain that produced them. trtexec
\texttt{-\/-layerPrecisions=l16-19:fp32} is \emph{effectively a no-op}
on explicit Q/DQ ONNX because the Q/DQ nodes themselves enforce the INT8
boundary at the tensor edges adjacent to blocks 16--19; whether the
kernel internal computation runs in FP32 or INT8 is invisible at the
tensor-precision interface. This is a non-obvious behavior of TensorRT
10.13 not surfaced explicitly in current vendor documentation; engineers
attempting mixed-precision via \texttt{-\/-layerPrecisions} should be
aware that the explicit Q/DQ surrounding block 16 input and block 19
output constitute the actual precision constraint, not the internal
kernel choice.

\textbf{Second, ModelOpt\textquotesingle s \texttt{disable\_quantizer}
and ONNX-level Q/DQ stripping are functionally equivalent.} Both prevent
Q/DQ insertion in the targeted block range, with the only difference
being the implementation layer (PyTorch calibration time vs ONNX library
post-export). The numerical equivalence (cos\_min 0.971 vs 0.9705,
speedup 2.41× vs 2.39×) confirms that the runtime path is determined by
the resulting ONNX graph topology, not the path that produced it.

\textbf{Third, cross-tool-chain validation should be standard practice
for negative-result claims.} A single-tool-chain mixed-precision study
reporting "ModelOpt + skip 16-19 fails to enter G2" leaves open whether
the failure is fundamental or implementation-specific. The convergence
across three independent paths transforms an isolated negative
observation into a falsifiable claim about the binding constraint.

\subsubsection{5.3 Implications for V1.3 QAT Future
Work}\label{53-implications-for-v13-qat-future-work}

The root cause analysis of §5.1 and the convergence proof of §5.2
jointly identify QAT as the only remaining path to satisfy the G2 cosine
constraint. The §4.7 V1.0.3 throughput findings strengthen this
conclusion along an independent axis: the BF16 dense path that
constitutes our V1.0.1 main deliverable runs at 96--99\% SM utilization
across three resolution regimes (r224 b1 N=2, r336 b8 N=1, r518 b8 N=1)
and at 98.74\% SM utilization at the G1 production configuration that
achieves 720 qps end-to-end. The dense compute path is therefore not
merely close to its precision ceiling; it is also at its hardware
compute ceiling. Any further throughput gain on this hardware-engine
combination must come from reducing per-inference compute volume --- and
the only published technique that does so while preserving the cos ≥
0.99 constraint is QAT-trained INT8 / FP8 / 2:4-sparse, which replaces
dense BF16 matmuls with reduced-precision Tensor Core operations.

We have published a complete V1.3 design ADR (ADR-011) detailing the
proposed implementation: starting from the SmoothQuant α=0.8 PTQ
initialization (the strongest INT8 baseline established in §4.3), apply
ModelOpt\textquotesingle s QAT mode for 1--5 fine-tuning epochs on
ImageNet val 50K with a small learning rate, re-export to Q/DQ ONNX, and
rebuild the TensorRT engine. The expected outcome is feat\_layer\_20
cos\_min ≥ 0.99 with speedup ≥ 3.0× (reflecting QAT\textquotesingle s
preservation of the underlying graph topology and thus the SmoothQuant
kernel selections), making this the first INT8 candidate inside the G2
ideal region. The V1.0.3 §4.7 SM-saturation evidence further suggests
that an INT8 QAT engine --- by reducing per-inference compute volume ---
can also lift the throughput ceiling above the dense BF16 720 qps
plateau, providing simultaneous gain on the precision and runtime axes.

We deliberately defer V1.3 implementation pending four launch
conditions: ImageNet val 50K access (formerly blocked by HF gated repo
\texttt{403\ GatedRepoError}; now resolved via Kaggle KGAT alternative
path), training resources (≥ 5 GPU-day on A100/H100 or ≥ 1 GPU-week on
RTX 5080), 1--2 month dedicated time including paper writing, and a
downstream task baseline (depth or segmentation FP32 + DPT head training
pipeline). All four conditions must be simultaneously met to ensure that
QAT is implemented under the right experimental controls; meeting only
some conditions risks producing a positive cosine result whose
generalization to downstream task quality cannot be verified.

\subsubsection{5.4 Methodological
Innovations}\label{54-methodological-innovations}

We identify three methodological patterns we believe generalize beyond
this specific project.

\textbf{Pure-Python testing for native tool-chain operations.} Splitting
graph manipulation into (a) pure-Python data-projection helpers that
operate on simple tuples and (b) thin remote-only driver scripts for
actual graph mutation enables full unit testing on developer
workstations without GPU or native ONNX dependencies. Our 271 pytest
cases across 111 source files run on macOS while the actual graph
rewriting runs on the Windows + RTX 5080 deployment node, eliminating
the developer-deployment friction that often plagues inference systems
research code.

\textbf{Bidirectional remote-sync separating binary from text-only
artifacts.} The \texttt{-\/-pull-reports} flag in our sync tool
reverse-pulls only text-only artifacts
(\texttt{.json/.md/.csv/.svg/.png/.jpg/.log/.txt}) via PowerShell-packed
zip archives, sidestepping flaky scp behavior on the cpolar SSH tunnel.
This separates large binary artifacts (ONNX, engines, weights) from
text-only reproducible artifacts in the round-trip, simplifying
day-to-day workflows without compromising reproducibility.

\textbf{Atomic SHA256 manifest with self-exclusion.} The
"manifest-generates-itself" 0-byte file pre-creation problem---where
shell-redirected manifest generation produces a 0-byte target file that
the manifest scanner records as a stale entry---is resolved by
\texttt{check\_assets.py\ -\/-output\ PATH}\textquotesingle s atomic
write (tempfile + \texttt{os.replace}) plus self-exclusion of the target
path. The pattern generalizes to any manifest-creates-itself situation
in reproducibility tooling.

\subsubsection{5.5 Comparison to Related
Work}\label{55-comparison-to-related-work}

Our findings extend three lines of prior work.

\textbf{ViT INT8 PTQ literature.} Standard ViT INT8 PTQ studies
{[}Frantar et al., 2023; Xiao et al., 2023; Liu et al., 2024{]} report
cos ≥ 0.95 thresholds suitable for image classification. Our cos ≥ 0.99
stringent regime --- motivated by downstream dense prediction
sensitivity --- exposes a regime where standard PTQ is insufficient on
ViT-L/16. We empirically validate that SmoothQuant\textquotesingle s
activation-only smoothing improves cos by \textasciitilde0.014 over
default PTQ on this model (0.968 vs 0.954, derived from §4.3
12-candidate scatter) but does not reach the 0.99 threshold, consistent
with SmoothQuant\textquotesingle s design as a smoothing-not-fine-tuning
method.

\textbf{DPT-style multi-scale fusion literature.} DPT {[}Ranftl et al.,
2021{]} recommends \texttt{{[}5,\ 11,\ 17,\ 23{]}} for ViT-L/24-block
backbones based on equally-spaced sampling. Our 4-layer ablation (§4.5)
empirically validates this layout achieves the highest inter-output
diversity but identifies a previously unreported 31.9× magnitude max/min
ratio that requires a fusion head capable of suppressing dominant
features. Our project\textquotesingle s \texttt{{[}4,\ 12,\ 16,\ 20{]}}
selection achieves a 2.5× tighter magnitude balance (12.6×) at a 22\%
diversity cost, which we interpret as a more practical default for
fusion-head architectures without explicit magnitude normalization.

\subsection{\texorpdfstring{\textbf{TensorRT mixed-precision
literature.} Existing TensorRT inference benchmarks for ViT models
{[}NVIDIA TensorRT release notes 10.0--10.13{]} do not report Blackwell
sm\_120 + ViT-L + cos ≥ 0.99 + cross-tool-chain comparison. Our work
documents the BF16 + Q/DQ Myelin Fill incompatibility (RoPE Constant
fails to lower under \texttt{-\/-bf16\ -\/-int8}) as a previously
undocumented intersection between TRT 10.13 and ViT-L deployment, and
quantifies the equivalence of \texttt{-\/-layerPrecisions} to
\texttt{disable\_quantizer} for explicit Q/DQ workflows. These findings
inform engineering decisions for practitioners deploying ViT-L on
Blackwell hardware in the TRT 10.13 release
window.}{TensorRT mixed-precision literature. Existing TensorRT inference benchmarks for ViT models {[}NVIDIA TensorRT release notes 10.0--10.13{]} do not report Blackwell sm\_120 + ViT-L + cos ≥ 0.99 + cross-tool-chain comparison. Our work documents the BF16 + Q/DQ Myelin Fill incompatibility (RoPE Constant fails to lower under -\/-bf16 -\/-int8) as a previously undocumented intersection between TRT 10.13 and ViT-L deployment, and quantifies the equivalence of -\/-layerPrecisions to disable\_quantizer for explicit Q/DQ workflows. These findings inform engineering decisions for practitioners deploying ViT-L on Blackwell hardware in the TRT 10.13 release window.}}\label{tensorrt-mixed-precision-literature-existing-tensorrt-inference-benchmarks-for-vit-models-nvidia-tensorrt-release-notes-1001013-do-not-report-blackwell-sm_120--vit-l--cos--099--cross-tool-chain-comparison-our-work-documents-the-bf16--qdq-myelin-fill-incompatibility-rope-constant-fails-to-lower-under---bf16---int8-as-a-previously-undocumented-intersection-between-trt-1013-and-vit-l-deployment-and-quantifies-the-equivalence-of---layerprecisions-to-disable_quantizer-for-explicit-qdq-workflows-these-findings-inform-engineering-decisions-for-practitioners-deploying-vit-l-on-blackwell-hardware-in-the-trt-1013-release-window}

\subsection{6. Limitations}\label{6-limitations}

We acknowledge four limitations that bound the generalizability of our
findings and motivate specific follow-up work.

\subsubsection{6.1 Dataset Proxy}\label{61-dataset-proxy}

ImageNet val 50K, the standard validation set used by most ViT
quantization studies, is gated under HuggingFace\textquotesingle s
\texttt{ILSVRC/imagenet-1k} repository and was inaccessible during the
project window (\texttt{403\ GatedRepoError}, no available proxy on the
GPU host network). We use Imagenette2-320 (10 classes, 13,000 images) as
a proxy distribution for both PTQ calibration and cosine evaluation.
While Imagenette is a strict subset of ImageNet\textquotesingle s 1000
classes, the smaller class diversity may underestimate worst-case PTQ
error in two ways: (a) under-represented texture / pose / occlusion
modes that exist in the full ImageNet distribution would not stress the
quantized network, and (b) the 10-class label structure may correlate
with feature regions that are easier to quantize than the full
1000-class manifold. We mitigate this by sampling 1000 images mutually
exclusive from the 500-image SmoothQuant calibration set, providing a
one-command swap-in path
(\texttt{scripts/export\_hf\_imagenet\_parquet\_images.py}) once
authorization is obtained, and acknowledging that the reported BF16
cos\_min ≥ 0.998 may be a slight overestimate when re-evaluated on full
ImageNet.

\subsubsection{6.2 Single-Hardware
Coverage}\label{62-single-hardware-coverage}

All measurements are taken on a single RTX 5080 (Blackwell sm\_120, 16
GB VRAM, 300 W TDP). We have not validated our findings on Ada Lovelace
(sm\_89), Hopper (sm\_90), or earlier consumer Blackwell SKUs. Three
findings are known or suspected to be hardware-specific: (a) FP8
hardware support requires Blackwell or Hopper---Ada
Lovelace\textquotesingle s lack of FP8 hardware would likely invert the
FP8 vs INT8 trade-off observed here; (b) the 16 GB VRAM constraint at
r518 batch ≥ 8 forces a profile narrowing that may inflate
per-(resolution, batch) timing variance compared to a 24+ GB device; (c)
the locked-clock 2752 MHz value is RTX 5080-specific and does not
transfer to other Blackwell SKUs with different boost clock targets. We
expect the qualitative findings (BF16 prefer dominance, three-tool-chain
mixed-precision convergence, upstream noise binding constraint) to
generalize but acknowledge that quantitative speedup numbers are
SKU-specific.

\subsubsection{6.3 TensorRT Version
Dependency}\label{63-tensorrt-version-dependency}

Several findings are specific to TensorRT 10.13.2.6: (a) the BF16 + Q/DQ
Myelin Fill incompatibility at the RoPE Constant node (failing with
"type assertion" under \texttt{-\/-bf16\ -\/-int8}); (b) the
\texttt{-\/-layerPrecisions=l16-19:fp32} no-op behavior on explicit Q/DQ
ONNX (the per-layer override is silently overridden by the surrounding
Q/DQ nodes); (c) the kernel selection for SmoothQuant α=0.8 INT8 engines
that yields the 3.48× speedup. Future TensorRT releases may close some
of these gaps---in particular, the RoPE Constant Myelin Fill issue is
plausibly a Blackwell-launch-window incompatibility that NVIDIA may
resolve. We have not exhaustively tested TRT 10.16.1 (the V1.0.1 ADR-008
ideal target) which may exhibit different behavior on the same engines.

\subsubsection{6.4 QAT Implementation
Deferred}\label{64-qat-implementation-deferred}

The V1.3 QAT path is provided as a complete design ADR (ADR-011) with
four launch conditions: ImageNet unblock, ≥ 5 GPU-day on A100/H100 (or ≥
1 GPU-week on RTX 5080), 1--2 month dedicated time including paper
writing, and a downstream task baseline (depth or segmentation FP32 +
DPT head training pipeline). None of the four are currently met. We
provide the design ADR but explicitly defer implementation, on the
principle that QAT under unmet preconditions risks producing a positive
cosine result whose generalization to downstream task quality cannot be
verified against a baseline.

\subsection{7. Conclusion}\label{7-conclusion}

This paper empirically maps the post-training quantization (PTQ)
operating region of TensorRT 10.13 inference for DINOv3 ViT-L/16
LVD-1689M on RTX 5080 (Blackwell sm\_120). Across 12 inference precision
candidates and three resolutions \{224, 336, 518\} on 1000 real images,
\textbf{BF16 prefer is the only candidate inside the (cos ≥ 0.99 ∧
speedup ≥ 2.2×) target region}, achieving 3.86× peak trtexec GPU-compute
speedup at r518 batch 8 with feat\_layer\_20 cos\_min ≥ 0.998 across all
three resolutions. The C++ runtime end-to-end peak speedup of 3.40×
confirms production deployability, and Python ↔ C++ runtime parity holds
bit-identically across all three resolutions at batch 1.

The central methodological contribution is a \textbf{three-tool-chain
convergence proof} for mixed-precision recovery: PyTorch
ModelOpt\textquotesingle s \texttt{disable\_quantizer},
TensorRT\textquotesingle s \texttt{-\/-layerPrecisions}, and ONNX-level
Q/DQ stripping all produce numerically equivalent results when forcing
transformer blocks 16--19 to non-INT8 precision (cos\_min span 0.0027,
b8 speedup span 0.04× across implementations that successfully force
FP32 fallback). This convergence eliminates "wrong tool chain" as an
explanation for the failure to enter the G2 ideal region and pinpoints
\textbf{upstream cumulative INT8 quantization noise from blocks 0--15}
as the binding constraint.

A V1.0.3 throughput-oriented serving extension (§4.7) further
establishes that the BF16 dense path is also at its hardware compute
ceiling: a multi-context inference pool achieves \textbf{720 qps}
end-to-end at r224 b=8 N=4 slots × 16 caller threads (2.094× over the
V1.0.2 single-stream baseline; \textbf{5.24× over the original FP32
trtexec baseline of 137 qps}) while driving SM utilization to 98.74\%
mean and 100\% p95. Concurrent execution under N=2 / N=4 produces output
bit-identical to the single-context reference (max\_abs\_error = 0,
cosine = 1.000). Three of three tested resolutions saturate at SM
88--99\% under their optimal batch size; r336 and r518 leave essentially
no idle GPU compute for additional contexts to exploit.

The constraint identification implies that quantization-aware
fine-tuning (QAT) is the only remaining path on \textbf{both} axes ---
the precision axis (escaping the upstream PTQ noise wall to enter the
cos ≥ 0.99 ∧ speedup ≥ 2.2× region) and the throughput axis (reducing
per-inference compute volume to lift the 720 qps in-process ceiling). We
publish a complete V1.3 QAT design ADR (ADR-011) with a four-condition
launch threshold and defer implementation pending all four conditions
being met.

We also contribute four methodological patterns we believe generalize
beyond this study: (a) splitting ONNX graph manipulation into
pure-Python identification helpers and thin remote-only mutation
drivers, enabling GPU-free unit testing of inference systems code; (b)
bidirectional remote-sync separating large binary artifacts (kept on the
GPU host) from text-only reproducible artifacts (round-tripped to
development hosts); (c) atomic SHA256 manifest with self-exclusion
eliminating the "manifest-generates-itself" 0-byte file pre-creation
problem; and (d) per-call SM utilization sampling at 100 ms cadence with
timing-coupled benchmark wallclock, enabling "throughput ceiling vs
hardware ceiling" decomposition analysis without specialized profiling
tools.

\subsection{\texorpdfstring{The complete reproducibility
kit---comprising a 56-row machine-readable benchmark matrix, twelve
V1.0.3 throughput data points across three resolutions × multiple batch
sizes × multiple pool capacities, eight SVG figures with a unified
\texttt{figures\_index.json} regeneration entry point, 271+ unit tests
across 111 source files, and an atomic SHA256 manifest covering 419+
artifact files---is released alongside this paper to support community
verification and
extension.}{The complete reproducibility kit---comprising a 56-row machine-readable benchmark matrix, twelve V1.0.3 throughput data points across three resolutions × multiple batch sizes × multiple pool capacities, eight SVG figures with a unified figures\_index.json regeneration entry point, 271+ unit tests across 111 source files, and an atomic SHA256 manifest covering 419+ artifact files---is released alongside this paper to support community verification and extension.}}\label{the-complete-reproducibility-kitcomprising-a-56-row-machine-readable-benchmark-matrix-twelve-v103-throughput-data-points-across-three-resolutions--multiple-batch-sizes--multiple-pool-capacities-eight-svg-figures-with-a-unified-figures_indexjson-regeneration-entry-point-271-unit-tests-across-111-source-files-and-an-atomic-sha256-manifest-covering-419-artifact-filesis-released-alongside-this-paper-to-support-community-verification-and-extension}

\subsection{8. References}\label{8-references}

{[}Numbered preliminary citation list; substitute with venue-specific
format during final submission. ACM Reference Format shown.{]}

{[}1{]} Maxime Oquab, Timothée Darcet, Théo Moutakanni, et al. 2024.
\textbf{DINOv2: Learning Robust Visual Features without Supervision.}
\emph{Transactions on Machine Learning Research} (TMLR).

{[}2{]} Meta AI Research. 2024. \textbf{DINOv3: Visual Self-Supervised
Learning at Scale.} \emph{(Place-holder for actual DINOv3 publication.)}

{[}3{]} René Ranftl, Alexey Bochkovskiy, and Vladlen Koltun. 2021.
\textbf{Vision Transformers for Dense Prediction.} In \emph{Proceedings
of the IEEE/CVF International Conference on Computer Vision (ICCV)}.
12179--12188.

{[}4{]} Guangxuan Xiao, Ji Lin, Mickael Seznec, Hao Wu, Julien Demouth,
and Song Han. 2023. \textbf{SmoothQuant: Accurate and Efficient
Post-Training Quantization for Large Language Models.} In
\emph{Proceedings of the 40th International Conference on Machine
Learning (ICML)}.

{[}5{]} Raghuraman Krishnamoorthi. 2018. \textbf{Quantizing deep
convolutional networks for efficient inference: A whitepaper.}
\emph{arXiv preprint arXiv:1806.08342}.

{[}6{]} Markus Nagel, Marios Fournarakis, Rana Ali Amjad, Yelysei
Bondarenko, Mart van Baalen, and Tijmen Blankevoort. 2021. \textbf{A
White Paper on Neural Network Quantization.} \emph{arXiv preprint
arXiv:2106.08295}.

{[}7{]} Zhihang Liu, et al. 2024. \textbf{FP8 Quantization for Vision
Transformers.} \emph{(Place-holder.)}

{[}8{]} NVIDIA Corporation. 2026. \textbf{TensorRT 10.13 Developer
Guide.} NVIDIA documentation.

{[}9{]} NVIDIA Corporation. 2026. \textbf{TensorRT Model Optimizer
Documentation.} NVIDIA documentation.

{[}10{]} NVIDIA Corporation. 2024. \textbf{NVIDIA Blackwell Architecture
White Paper.}

{[}11{]} Hugo Touvron, Matthieu Cord, and Hervé Jégou. 2022.
\textbf{DeiT III: Revenge of the ViT.} In \emph{Proceedings of the
European Conference on Computer Vision (ECCV)}.

{[}12{]} Elias Frantar, Saleh Ashkboos, Torsten Hoefler, and Dan
Alistarh. 2023. \textbf{GPTQ: Accurate Post-Training Quantization for
Generative Pre-trained Transformers.} In \emph{Proceedings of the
International Conference on Learning Representations (ICLR)}.

{[}13{]} NVIDIA Corporation. 2024. \textbf{NVIDIA Triton Inference
Server User Guide.} NVIDIA documentation.
\url{https://docs.nvidia.com/deeplearning/triton-inference-server/}

{[}14{]} NVIDIA Corporation. 2024. \textbf{CUDA Multi-Process Service
(MPS) Documentation.} \url{https://docs.nvidia.com/deploy/mps/}

{[}15{]} Asit Mishra, Jorge Albericio Latorre, Jeff Pool, Darko Stosic,
Dusan Stosic, Ganesh Venkatesh, Chong Yu, and Paulius Micikevicius.
2021. \textbf{Accelerating Sparse Deep Neural Networks.} \emph{arXiv
preprint arXiv:2104.08378}.

{[}16{]} Maxim Milakov and Natalia Gimelshein. 2018. \textbf{Online
Normalizer Calculation for Softmax.} \emph{arXiv preprint
arXiv:1805.02867}. \emph{(Cited for Tensor Core BF16 dense path
utilization-ceiling structural bound discussed in §4.7.)}

\begin{center}\rule{0.5\linewidth}{0.5pt}\end{center}

\subsection{Acknowledgments}\label{acknowledgments}

This work uses Meta DINOv3 ViT-L/16 LVD-1689M
(\texttt{facebook/dinov3-vitl16-pretrain-lvd1689m}) under the DINOv3
License. Built with DINOv3. The repository copy of the license is
available at \texttt{LICENSES/DINOv3\_LICENSE.md}. NVIDIA TensorRT
10.13.2.6 and Model Optimizer are used per their respective licenses.

{[}Funding sources, advisor acknowledgments, peer feedback per
submission requirements.{]}

\begin{center}\rule{0.5\linewidth}{0.5pt}\end{center}

\subsection{Word Count Summary}\label{word-count-summary}

{\def\LTcaptype{none} % do not increment counter
\begin{longtable}[]{@{}lrrr@{}}
\toprule\noalign{}
Section & Word count V1.0.0 & V1.0.1 delta & V1.0.1 total \\
\midrule\noalign{}
\endhead
\bottomrule\noalign{}
\endlastfoot
Abstract (EN) & 312 & +168 & 480 \\
Abstract (中文) & 720 chars & +360 chars & 1,080 chars \\
§ 1 Introduction & 1,878 & 0 & 1,878 \\
§ 2 Literature Review & 1,194 & 0 & 1,194 \\
§ 3 Methodology & 1,703 & 0 & 1,703 \\
§ 4 Results & 1,633 & +1,150 (§4.7) & 2,783 \\
§ 5 Discussion & 1,314 & +180 (§5.3) & 1,494 \\
§ 6 Limitations + § 7 Conclusion & 885 & +220 (§7) & 1,105 \\
\textbf{Total English body} & \textbf{\textasciitilde8,919 words} &
\textbf{+1,718 words} & \textbf{\textasciitilde10,637 words} \\
\end{longtable}
}

V1.0.0 fits the 6,000-8,000 word workshop target. V1.0.1 (with V1.0.3
§4.7 throughput-saturation extension) fits the 10,000-12,000 word full
conference target after light editing. Eligible venue formats include
MLSys (≤ 9 pages double-column), NeurIPS Datasets \& Benchmarks (≤ 9
pages), and ICLR Tiny Papers / ML for Systems workshops.

\end{document}
